\chapter{道具の準備：Python・NumPy・PyTorch のテンソル}\label{ch:pytorch}
\bookmap{1}

\begin{chapterlead}
\textbf{ここまで}: \chapref{ch:math}で、ベクトル・行列・微分・確率という、数式の「言葉」をそろえました。\\[1mm]
\textbf{この章の問い}: 数式どおりの計算を、何百万個もの数に対して、速く・短く・間違えずに実行するには、どんな道具をどう使えばよいのでしょうか。\\[1mm]
\textbf{この章を終えると}:
\begin{itemize}
\item テンソルの「形」を読み、変え、形の違う相手との計算を \code{for} 文なしで書けます。
\item 数式 $\mW\vx+\vb$ を多数のデータに一度に適用するコードが読め、自動微分で勾配を求められます。
\item softmax と交差エントロピーを自分で書き、本物のモデルの損失と一致することを確かめられます。
\end{itemize}
\end{chapterlead}

\begin{readingguide}
\item[\rgpre] この章は、通読する章ではなく、必要な節を読んで、あとは辞書として戻る章です。\chapref{ch:math}の行列積と勾配を思い出せれば読めます。
\item[\rgmust] \secref{sec:pt-tools}と、\secref{sec:pt-tensor}〜\secref{sec:pt-autograd}：形、ブロードキャスト、行列積 \code{X @ W.T + b}、自動微分。後の章のコードは、この四つで書かれています。章末の演習問題は、初読では最初の三つで十分です。
\item[\rgskim] \secref{sec:pt-speed}：「数を 1 個ずつ回る \code{for} 文は、配列演算より 1,000 倍以上遅い」という結論だけで先へ進めます。\secref{sec:pt-tensor}の「型と、テンソルの作り方」と、\secref{sec:pt-final}の腕試しも、札に書いた結論を押さえれば足ります。
\item[\rglater] \secref{sec:pt-reshape}の最後の一覧（\tabref{tab:pt-shape-ops}）と、\secref{sec:pt-final}の前半（エラー文の読み方）：使う場面が来たら戻ります。書き方を忘れたときは、章末の早見表（\tabref{tab:pt-cheatsheet}）から引けます。
\end{readingguide}

\chapref{ch:intro}では、できあがったモデルに「どうも仕方が」の続きを予測させました。
あの 1 回の予測の裏では、約 260 万個の数を使った計算が走っています。
中身は、\chapref{ch:math}で学んだ内積・行列積・softmax の組み合わせです。
この計算を、Python の \code{for} 文で数を 1 個ずつ処理して実行すると、遅すぎて使いものになりません。
基盤モデルの学習に出てくる行列積のうち、たった 1 種類を \code{for} 文で実行するだけで、約 8 日かかります。
同じ計算が、PyTorch（パイトーチ）という道具を使えば、約 8 分で終わります。
PyTorch の中心には、数をまとめて入れておく箱である「テンソル」があります。
この章の合言葉は「\textbf{形を見よ}」です。
テンソルの計算で起きる間違いの大半は、数の並び方、つまり形の取り違えだからです。

\section{道具をそろえる}\label{sec:pt-tools}

最初に、道具を手元に用意します。何を入れ、それぞれは何の役をするのでしょうか。
本書で使う道具は三つです。
Python は、計算の段取りを書くプログラミング言語です。
NumPy（ナムパイ）は、数の並びをまとめて計算するためのライブラリです~\cite{harris2020numpy}。
数を並べたものを、プログラミングでは\term{配列}{array}と呼びます。
PyTorch\cite{paszke2019pytorch} は、NumPy とよく似た書き方で配列を扱えるライブラリです。
そのうえ、NumPy にない機能を二つ持っています。
微分を自動で計算する機能（\secref{sec:pt-autograd}）と、GPU という計算装置で大量の計算を速く実行する機能（\chapref{ch:gpu}）です。
どちらも学習に欠かせないので、本書の主役は PyTorch です。

Python（3.11 以降）の入った PC で \code{pip install torch numpy} を実行すれば、準備は終わりです。
詳しい手順は\appref{app:setup}にまとめてあります。本書のコードは、GPU がなくても動きます。
さっそく動かしてみます。

\snippet{最初のテンソル}{code:pt-first}{ch02/hello_tensor.py}{first}

\outputfile{ch02/out/hello_tensor_1.txt}

1 行目の \code{import torch} は、PyTorch を使えるようにする宣言です。この章の以後のコードでは、この行を省いて載せます。
3〜4 行目の \code{torch.tensor} は、Python のリストからテンソルを作る関数です。
ここでは、$\vx=(1,2,3)$ と $\vy=(10,20,30)$ という二つのベクトルを作りました。
出力の 1 行目が \code{x + y} の結果で、中身は $(1+10,\,2+20,\,3+30)$ です。
\code{11.} のように数の後ろに付いた点は、「小数として持っている」という印です。
2 行目はスカラー倍 $2\vx$、3 行目は \code{x @ y} で、内積 $1\times 10+2\times 20+3\times 30=140$ です（数 1 個でも、結果はテンソルです）。
記号 \code{@} は、内積と行列積を表す演算子です（\secref{sec:pt-matmul}）。
数式の $\vx+\vy$、$2\vx$、$\vx^\T\vy$ が、ほぼそのままの見た目でコードになりました。\code{for} 文は一つも書いていません。

なお、NumPy の配列とテンソルは、\code{torch.from_numpy(a)} と \code{t.numpy()} で行き来できます（\file{ch02/numpy_bridge.py}）。
テンソルを置いて計算する装置（CPU か GPU）を\term{デバイス}{device}と呼びますが、\chapref{ch:gpu}までは意識しなくてかまいません。

この節では、道具をそろえ、数式に近い見た目で計算が書けることを確かめました。

\section{なぜ for 文ではなく配列演算か}\label{sec:pt-speed}
\skimread{結論は「数を 1 個ずつ回る \code{for} 文は、配列演算より 1,000 倍以上遅い。だから、数はテンソルにまとめて計算する」です。測り方と、遅くなる理由は、2 周目に読めば足ります。}

前節では、内積を \code{x @ y} の 1 行で計算しました。
しかし、すでに知っている \code{for} 文でも、同じ計算は書けます。
わざわざテンソルを使う理由は何でしょうか。
答えは「速さ」です。どのくらい違うのかを、行列積で測ります。

行列積は、$C_{ij}=\sum_{k}A_{ik}B_{kj}$ で定義される計算でした（\chapref{ch:math}）。
この式を Python の \code{for} 文でそのまま書くと、$i$、$j$、$k$ についての 3 重の \code{for} 文になります。
行列を「リストのリスト」で表し、いちばん内側で \code{s += A[i][k] * B[k][j]} を実行します（\file{ch02/speed.py} の関数 \code{matmul_loop}）。
$n$ 行 $n$ 列どうしなら、この掛け算を $n\times n\times n=n^{3}$ 回実行します。$i$、$j$、$k$ が、それぞれ $n$ 通りの値を取るからです。
$n=256$ なら、$256^{3}$ で約 1,678 万回です。

同じ計算は、NumPy でも PyTorch でも、\code{A @ B} の 1 行です。
乱数で埋めた $n$ 行 $n$ 列の行列を二つ用意し、三つの方法で積を計算して、時間を測りました。
\figref{fig:pt-speed}が結果です。
縦軸は 1 目盛りで 10 倍になる対数目盛なので、線の上下の間隔が「何倍遅いか」を表します。

\begin{figure}[tbp]
\centering
\begin{tikzpicture}
\begin{axis}[xmode=log,ymode=log,log basis x=2,width=0.72\linewidth,height=0.31\linewidth,
  xlabel={行列の大きさ $n$（$n$ 行 $n$ 列どうしの積）},ylabel={時間 [ミリ秒]},
  xtick={8,16,32,64,128,256},xticklabels={8,16,32,64,128,256},
  legend pos=north west,legend cell align=left,legend style={font=\footnotesize}]
\addplot[bookaccent,mark=*,mark size=1.6pt] table[col sep=comma,x=n,y=loop_ms]{code/ch02/out/speed_matmul.csv};
\addlegendentry{\texttt{for} 文}
\addplot[bookgreen,mark=square*,mark size=1.6pt] table[col sep=comma,x=n,y=numpy_ms]{code/ch02/out/speed_matmul.csv};
\addlegendentry{NumPy}
\addplot[bookblue,mark=triangle*,mark size=2pt,dashed] table[col sep=comma,x=n,y=torch_ms]{code/ch02/out/speed_matmul.csv};
\addlegendentry{PyTorch}
\end{axis}
\end{tikzpicture}
\caption{行列積にかかる時間の実測（両軸とも対数目盛）。数回測って最も速かった回の値である。条件をそろえるため、どの方法にも CPU の計算回路（コア）を 1 個だけ使わせた。}
\label{fig:pt-speed}
\end{figure}

右端の $n=256$ を見てください。
\code{for} 文は 857 ミリ秒、PyTorch は 0.65 ミリ秒でした。
答えは同じなのに、\code{for} 文は 1,317 倍の時間を使っています（倍率は、丸める前の測定値から求めました）。
二つの方法の答えを比べると、差は最大でも約 $10^{-5}$ でした。小数の丸め方の違いによるもので、同じ答えとみなせます。
NumPy と PyTorch の線は、$n$ が大きいところでほぼ重なっています。中で、同じ種類の高速なプログラムを呼んでいるからです。
差は $n$ が大きいほど開きます（$n=8$ では 7 倍です）。モデルが扱う行列は大きいので、差が大きい側が現実です。

なぜ、これほど違うのでしょうか。
$n^{3}$ 回の掛け算は、どの方法でも必要です。違うのは、\textbf{1 回あたりの速さ}です。理由は二つあります。

\begin{figure}[tbp]
\centering
\begin{tikzpicture}[x=1mm,y=1mm,font=\footnotesize]
  \node[anchor=west] at (-2,25) {\textbf{Python のリスト}\ \ \texttt{[1.0, 2.0, 3.0]}};
  \foreach \i in {0,1,2}
    \node[draw=black!45,fill=figgray,minimum width=8mm,minimum height=5.5mm,inner sep=0pt] (p\i) at (4+8*\i,17) {$\bullet$};
  \node[lbl,anchor=west] at (25.5,17) {「置き場所」の列};
  \node[plainbox,minimum width=14mm,font=\scriptsize] (o0) at (3,3) {型: 小数\\値: 1.0};
  \node[plainbox,minimum width=14mm,font=\scriptsize] (o2) at (24,5.5) {型: 小数\\値: 3.0};
  \node[plainbox,minimum width=14mm,font=\scriptsize] (o1) at (45,2) {型: 小数\\値: 2.0};
  \draw[arr] (p0.south) -- (o0.north);
  \draw[arr] (p1.south) to[out=-70,in=150] (o1.north west);
  \draw[arr] (p2.south) -- (o2.north);
  \begin{scope}[shift={(72,0)}]
    \node[anchor=west] at (-2,25) {\textbf{配列・テンソル}};
    \node[graybox,minimum width=30mm,font=\scriptsize] (h) at (15,17) {形: (3,)\quad 型: 小数（全要素共通）};
    \foreach \i/\v in {0/1.0,1/2.0,2/3.0}
      \node[draw=black!45,fill=figblue,minimum width=10mm,minimum height=5.5mm,inner sep=0pt] (c\i) at (5+10*\i,4) {\v};
    \draw[arr] (h.south) -- (c1.north);
  \end{scope}
\end{tikzpicture}
\caption{リストと配列の、メモリ上の姿。リスト（左）は数の「置き場所」を並べ、数の本体は別々の場所にある。
配列（右）は、同じ型の数そのものを一続きに並べる。}
\label{fig:pt-memory}
\end{figure}

一つ目は、\textbf{数の持ち方}です（\figref{fig:pt-memory}）。
Python のリストは、どんな種類のものでも入れられる代わりに、数そのものではなく、数の「置き場所」を並べています。
掛け算を 1 回するたびに、置き場所をたどり、型を調べ、結果を入れる新しい場所を用意します。
配列は、全部の要素が同じ型だと最初に決めてしまい、数そのものを一続きに並べるので、どちらの手間もありません。

二つ目は、\textbf{計算の実行のされ方}です。
\code{A @ B} と書くと、Python は「この二つの行列の積を求めよ」という一つの命令だけを出します。
あとの $n^{3}$ 回の掛け算は、あらかじめ機械向けに翻訳され、磨き上げられたプログラムが一気に実行します。
100 個の荷物を 1 個ずつ伝票を書いて運ぶか、トラックに積んで 1 回で運ぶか、の違いです。
要素を一つずつ処理する \code{for} 文を、配列全体への演算に書き直すことを、\term{ベクトル化}{vectorization}と呼びます。

最後に、本書の基盤モデルに実際に出てくる大きさでも測りました。
128 行 192 列と 192 行 4,096 列の行列の積で、128 個のトークンのそれぞれに、4,096 個の候補のスコアを付ける計算です。
掛け算の回数は $128\times 192\times 4{,}096$ で、約 1 億回です。
結果は、\code{for} 文が 5.36 秒、PyTorch が 3.68 ミリ秒で、差は 1,455 倍でした（これも、丸める前の値の比です）。
基盤モデルの学習では、パラメータの書き換えを 4,000 回繰り返し、1 回ごとに 32 個の文を使います。この行列積の回数は、$4{,}000\times 32=128{,}000$ 回です。
\code{for} 文なら、$128{,}000\times 5.36=686{,}080$ 秒かかります。1 日は 86,400 秒なので、$686{,}080\div 86{,}400\approx 7.9$ で、\textbf{約 8 日}です。
PyTorch なら、$128{,}000\times 3.68$ ミリ秒で約 471 秒です。60 で割ると約 7.9 分で、\textbf{約 8 分}です。
なお、避けるべきなのは、\textbf{数の一つひとつを回る} \code{for} 文だけです。
「書き換えを 4,000 回繰り返す」といった、回数の少ない \code{for} 文は、これからも普通に使います。

この節では、同じ行列積でも、\code{for} 文と配列演算とで、速さが 1,000 倍以上違うことを実測しました。

\section{テンソル：形と型}\label{sec:pt-tensor}

配列演算が速いことは分かりました。
では、その配列、つまり PyTorch のテンソルとは、どんなものなのでしょうか。
\figref{fig:pt-memory}の右側のとおり、テンソルは、一続きに並んだ数と、「形」と「型」という二つの情報からできています。

\subsection*{形：どの向きに、いくつ並んでいるか}

\chapref{ch:math}では、数の並べ方を、スカラー・ベクトル・行列・テンソルと呼び分けました。
PyTorch では、この四つを区別せず、すべて\textbf{テンソル}と呼びます。違いは「形」だけです。
数が並ぶ向きの数を\jterm{階数}と呼び（コードでは \code{.ndim}）、順に 0 階、1 階、2 階、3 階のテンソルとも言います。

\snippet{いろいろな形のテンソル}{code:pt-dims}{ch02/shape_dtype.py}{dims}

\outputfile{ch02/out/shape_dtype_1.txt}

1〜4 行目で、数 1 個、ベクトル、行列を作っています。行列は「行のリスト」のリストとして書きます。
5 行目の \code{torch.zeros(2, 3, 4)} は、すべて 0 で埋めた、形 $(2,3,4)$ のテンソルです。
6 行目の \code{.shape} が、テンソルの\textbf{形}を返します。
出力の 1 行目の \code{torch.Size([2, 3])} は、2 行 3 列という意味です。本書では、これを $(2,3)$ と書きます。
2 行目以降は、四つのテンソルの形を、Python のタプルに直して表示したものです。
ベクトル \code{v} の形 \code{(3,)} は、「3 個の数が 1 列に並ぶ」という意味です（カンマは、要素が 1 個のタプルの書き方です）。
スカラー \code{s} の形は、空のタプル \code{()} です。並ぶ向きが一つもない、ということです。

この「向き」のことを\term{軸}{axis}と呼び、形の左から順に、軸 0、軸 1、軸 2 と番号を付けます（\figref{fig:pt-axes}）。
行列では、軸 0 が縦（行の番号が増える向き）、軸 1 が横（列の番号が増える向き）です。
形 $(2,3,4)$ のテンソルは、3 行 4 列の行列が 2 枚重なったもので、軸 0 は「何枚目か」を表します。
行と列は、いつも形の\textbf{右端の二つ}です。軸が増えるときは左に足されるので、縦は軸 1、横は軸 2 にずれます。
逆に、形 $(3,)$ のベクトルは、軸が一つあるだけで、縦か横かの区別を持ちません。
出力の右端の「要素数」（\code{.numel()}）は、形に並ぶ数をすべて掛けたものです。\code{T} なら $2\times 3\times 4=24$ 個です。

\begin{figure}[tbp]
\centering
\begin{tikzpicture}[x=5.5mm,y=5.5mm,font=\footnotesize]
  \fill[figgreen] (0,0) rectangle (3,2);
  \draw[black!45,step=1] (0,0) grid (3,2);
  \foreach \v/\x/\y in {1/0/1,2/1/1,3/2/1,4/0/0,5/1/0,6/2/0} \node at (\x+0.5,\y+0.5) {\v};
  \draw[arr,bookaccent] (-0.5,2) -- (-0.5,0) node[midway,left,text=bookaccent] {軸 0};
  \draw[arr,bookblue] (0,2.5) -- (3,2.5) node[midway,above,text=bookblue] {軸 1};
  \node[anchor=west,align=left] at (3.6,1) {\texttt{M}\ \ 形 $(2,3)$\\軸 0 の長さ 2\\軸 1 の長さ 3};
  \begin{scope}[shift={(12.5,-0.9)}]
    \foreach \k in {1,0}{
      \begin{scope}[shift={(0.8*\k,0.8*\k)}]
        \fill[figpurple] (0,0) rectangle (4,3);
        \draw[black!45,step=1] (0,0) grid (4,3);
      \end{scope}}
    \draw[arr,bookaccent] (-0.5,3) -- (-0.5,0) node[midway,left,text=bookaccent] {軸 1};
    \draw[arr,bookblue] (0,-0.5) -- (4,-0.5) node[midway,below,text=bookblue] {軸 2};
    \draw[arr,bookgreen] (4.3,0.1) -- (5.2,1.0) node[midway,below right,text=bookgreen] {軸 0（何枚目か）};
    \node[anchor=west,align=left] at (5.6,3.1) {\texttt{T}\ \ 形 $(2,3,4)$\\3 行 4 列が 2 枚};
  \end{scope}
\end{tikzpicture}
\caption{軸と形。形は、各軸の長さを左から順に並べたものである。行列（左）は軸を二つ、3 階のテンソル（右）は軸を三つ持つ。}
\label{fig:pt-axes}
\end{figure}

\subsection*{型と、テンソルの作り方}
\skimread{結論は三つです。小数は \code{float32}、番号は \code{int64} という型で持ちます。型の違う相手とは、行列積でエラーになります。乱数は、種を固定すれば毎回同じ数になります。テンソルを作る関数は、出てきたときにここへ戻れば足ります。}

テンソルのもう一つの情報が、\term{データ型}{dtype}です。型は \code{.dtype} で調べられます。本書で使う型は、\tabref{tab:pt-dtype}の三つでほぼすべてです。

\begin{table}[tbp]
\centering
\caption{本書で使う主なデータ型。名前の後ろの数字は、1 個の数に使うビット数である（8 ビットが 1 バイト）。}
\label{tab:pt-dtype}
\small
\begin{tabular}{@{}llrl@{}}
\toprule
データ型 & 中身 & 1 要素 & 主な用途 \\
\midrule
\texttt{torch.float32} & 小数（有効数字 約 7 桁） & 4 バイト & モデルのパラメータ、計算の途中の値 \\
\texttt{torch.int64} & 整数 & 8 バイト & トークンの番号、正解の番号 \\
\texttt{torch.bool} & \texttt{True} か \texttt{False} & 1 バイト & 条件に合うかどうかの印 \\
\bottomrule
\end{tabular}
\end{table}

\code{float32} は、32 ビット（4 バイト）で小数を表す形式で、\term{浮動小数点数}{floating-point number}と呼ばれます。
表せる数には限りがあり、正しいのはおよそ 7 桁目までです。
たとえば 0.1 を入れて取り出すと、\code{0.10000000149...} が返ってきます。
モデルの学習にはこの精度で十分なので、メモリの少ない \code{float32} が標準です。
型の違いは、エラーの原因になります。
\code{torch.tensor([1, 2, 3])} のように整数だけで作ると \code{int64} になり、\code{float32} との行列積ができません。
小数として使うなら、\code{[1.0, 2.0, 3.0]} と小数点を付けて書くか、\code{.float()} で変換します。

大きなテンソルは、リストでは書けないので、形を指定して作ります。
先ほどの \code{torch.zeros} のほか、\code{torch.arange(6)} は、0 から 5 までの整数を並べたテンソルを作ります。
\code{torch.randn(2, 3)} は、標準正規分布（\chapref{ch:math}）に従う\textbf{乱数}で埋めます。学習を始めるとき、モデルのパラメータは乱数で埋めます。
乱数は呼ぶたびに違う数になりますが、直前に \code{torch.manual_seed(0)} で\term{乱数の種}{seed}を決めると、毎回同じ数が出ます（\file{ch02/random_seed.py}）。
コンピュータの乱数は、種を出発点にして、決まった手順で次々に作られる数だからです。出発点が同じなら、出てくる数の並びも同じになります。
何度実行しても同じ結果になる性質を、\term{再現性}{reproducibility}と呼びます。本書のコードは、必ず種を固定しています。

\subsection*{モデルの正体はテンソルの集まり}

ここまでの知識で、本物のモデルの中をのぞけます。
\chapref{ch:intro}で動かした基盤モデルを読み込み、中にあるテンソルの名前・形・型を表示しました（\file{ch02/model_tensors.py}）。
先頭の 5 個と、全体の合計は、次のとおりです。

\outputfile{ch02/out/model_tensors.txt}

名前の意味は、まだ分からなくてかまいません（\chapref{ch:embedding}から\chapref{ch:transformer}で一つずつ作ります）。
見てほしいのは形と型です。
最初のテンソルは形 $(4096,192)$ の行列で、4,096 個のトークンのそれぞれに、192 個の数が割り当てられています。型はすべて \code{float32}（小数を 1 個あたり 4 バイトで持つ型）です。
テンソルは全部で 36 個、要素数の合計は 2,583,936 個でした。\chapref{ch:intro}で「約 260 万個のパラメータ」と呼んだものの内訳です。
出力の最後の行のメモリは、要素数に 4 バイトを掛けたものです。$2{,}583{,}936\times 4$ は約 1,034 万バイトで、約 10.3 MB です。
モデルとは、形の決まったテンソルの集まりで、学習とは、その中の数を少しずつ書き換えていく作業です。

\subsection*{要素を取り出す：インデックスとスライス}

テンソルの一部を取り出すには、角括弧を使います。
位置を番号で指定する操作を\term{インデックス}{indexing}、範囲で指定する操作を\term{スライス}{slicing}と呼びます。

\snippet{インデックスとスライス}{code:pt-indexing}{ch02/indexing.py}{basic}

\outputfile{ch02/out/indexing_1.txt}

行列 \code{A} は、十の位を行の番号、一の位を列の番号にしてあります。
4 行目の \code{A[1, 2]} は、行 1・列 2 の要素で、結果は 12 です。
\textbf{コードの番号は 0 から始まります}。
数式で 1 から数えて $A_{23}$ と書く要素が、コードでは \code{A[1, 2]} です。
1 個だけ取り出しても、結果は 0 階のテンソル \code{tensor(12)} です。Python の普通の数にするには、\code{.item()} を付けます。
5〜7 行目が、行・列・範囲の取り出しです（\figref{fig:pt-indexing}）。
「\code{:}」は「その軸は全部」、\code{0:2} は「0 以上 2 未満」を表します。終わりの番号は含みません。
8 行目のように、番号を並べたテンソルを渡すと、その番号の行が順に並びます。
この操作は、\chapref{ch:embedding}で「トークンの番号から、そのトークンの行を引く」ために使います。
9 行目の \code{A > 20} は、要素ごとに「20 より大きいか」を調べた、\code{True} と \code{False} のテンソルです（型は \code{bool}）。
これを角括弧に渡すと、\code{True} の位置の要素だけが、1 列に並んで返ります。

\begin{figure}[tbp]
\centering
\begin{tikzpicture}[x=5.2mm,y=5.2mm,font=\footnotesize]
\foreach \p/\ttl in {0/{\texttt{A[1]}\ \ 形 $(4,)$},1/{\texttt{A[:, 2]}\ \ 形 $(3,)$},2/{\texttt{A[0:2, 1:3]}\ \ 形 $(2,2)$}}{
  \begin{scope}[shift={(7.6*\p,0)}]
    \ifcase\p \fill[figorange] (0,1) rectangle (4,2);
    \or \fill[figorange] (2,0) rectangle (3,3);
    \or \fill[figorange] (1,1) rectangle (3,3);
    \fi
    \draw[black!45,step=1] (0,0) grid (4,3);
    \foreach \r/\y in {0/2,1/1,2/0}{
      \foreach \c in {0,1,2,3}{
        \pgfmathtruncatemacro{\v}{10*\r+\c}
        \node at (\c+0.5,\y+0.5) {\v};}
      \node[lbl] at (-0.4,\y+0.5) {\r};}
    \foreach \c in {0,1,2,3} \node[lbl] at (\c+0.5,3.35) {\c};
    \node at (2,-0.6) {\ttl};
  \end{scope}}
\end{tikzpicture}
\caption{インデックスとスライスで取り出される部分（色の付いたマス）。灰色の数字は、行と列の番号（0 始まり）である。}
\label{fig:pt-indexing}
\end{figure}

一つ、落とし穴があります。
番号や範囲で取り出した結果は、数のコピーではありません。元のテンソルのメモリの一部を、のぞき窓のように見せているだけです（\term{ビュー}{view}と呼びます）。
\code{row = A[0]} として \code{row[0] = 99} と書き換えると、元の \code{A} も変わります。コピーが欲しいときは、\code{A[0].clone()} とします（8〜9 行目の取り出し方は、コピーを作ります）。

この節では、テンソルが形と型を持つことを学びました。本物のモデルも、形の決まったテンソルの集まりでした。

\section{形を変える}\label{sec:pt-reshape}

テンソルの形を読めるようになりました。次は、形を\textbf{変える}操作です。
モデルの中では、同じ数の集まりを、別の形として読み直す場面がたくさんあります。
たとえば、「8 個の文 $\times$ 128 トークン」の予測を、損失の計算のために「1,024 個の予測」として 1 列に並べ直します。
形を変えるとき、中の数はどうなるのでしょうか。

\snippet{reshape で形を変える}{code:pt-reshape}{ch02/reshape.py}{reshape}

\outputfile{ch02/out/reshape_1.txt}

1 行目の \code{torch.arange(12)} は、0 から 11 までの整数を順に並べた、形 $(12,)$ のテンソルを作ります。
2 行目の \code{x.reshape(3, 4)} は、12 個の数を 3 行 4 列に並べ直します。
出力を見ると、0 から順に、4 個ずつ区切って行にしています。
\figref{fig:pt-reshape}が、この操作の正体です。
テンソルの数は、メモリの上では、形に関係なくいつも 1 列に並んでいます。
形とは、その 1 列を\textbf{どこで区切って読むか}という約束です。
\code{reshape} は、数の並びには触れず、区切り方だけを取り替えます。
区切り方を変えるだけなので、要素数は変えられません。
12 個の数は、3 行目のように $(2,6)$ にも $(2,2,3)$ にもできますが、$(5,3)$ にはできません。$2\times 6$ も $2\times 2\times 3$ も 12 ですが、$5\times 3$ は 15 で、12 と合わないからです。
4 行目の \code{-1} は、「この軸の長さは、残りから計算してほしい」という意味です。12 個を 4 行にするなら、列は 3 です。

\begin{figure}[tbp]
\centering
\begin{tikzpicture}[x=5.2mm,y=5.2mm,font=\footnotesize]
  \fill[figgray] (0,4.6) rectangle (12,5.6);
  \draw[black!45] (0,4.6) rectangle (12,5.6);
  \foreach \i in {1,...,11} \draw[black!45] (\i,4.6) -- (\i,5.6);
  \foreach \i in {0,...,11} \node at (\i+0.5,5.1) {\i};
  \node[lbl,anchor=west] at (12.3,5.1) {メモリの上の並び（1 列）};
  \begin{scope}[shift={(0.5,0)}]
    \fill[figblue] (0,2) rectangle (4,3); \fill[figgreen] (0,1) rectangle (4,2); \fill[figorange] (0,0) rectangle (4,1);
    \draw[black!45,step=1] (0,0) grid (4,3);
    \foreach \r/\y in {0/2,1/1,2/0}{\foreach \c in {0,1,2,3}{
      \pgfmathtruncatemacro{\v}{4*\r+\c}\node at (\c+0.5,\y+0.5) {\v};}}
    \node[anchor=west] at (4.2,0.4) {\texttt{x.reshape(3, 4)}};
  \end{scope}
  \begin{scope}[shift={(11.5,1)}]
    \fill[figblue] (0,1) rectangle (6,2); \fill[figgreen] (0,0) rectangle (6,1);
    \draw[black!45,step=1] (0,0) grid (6,2);
    \foreach \r/\y in {0/1,1/0}{\foreach \c in {0,...,5}{
      \pgfmathtruncatemacro{\v}{6*\r+\c}\node at (\c+0.5,\y+0.5) {\v};}}
    \node at (3,-0.6) {\texttt{x.reshape(2, 6)}};
  \end{scope}
  \draw[arr] (2.5,4.4) -- (2.5,3.2) node[midway,right,lbl] {4 個ずつ区切る};
  \draw[arr] (11,4.4) -- (13,3.2) node[midway,right,lbl] {\ 6 個ずつ区切る};
\end{tikzpicture}
\caption{形は「読み方」である。メモリの上の数の並び（上）は変わらない。
4 個ずつ区切って読めば形 $(3,4)$、6 個ずつなら形 $(2,6)$ になる。色は行を表す。}
\label{fig:pt-reshape}
\end{figure}

\subsection*{転置は reshape とは違う}

形を変える操作には、もう一つ、\textbf{転置}があります。
転置 $\mA^\T$ は、行と列を入れ替える操作でした（\chapref{ch:math}）。コードでは \code{A.T} と書きます。
2 行 3 列の行列を転置すると、3 行 2 列になります。
それなら、\code{A.reshape(3, 2)} と同じでしょうか。
\figref{fig:pt-transpose}が、実際に計算した結果です（\file{ch02/reshape.py}）。
形はどちらも $(3,2)$ ですが、\textbf{中身の並びが違います}。
転置では、元の行 $(0,1,2)$ が、そのまま縦の列になっています。区切り方ではなく、数を読む順番を変える操作です。
\code{reshape} は、0 から順に 2 個ずつ区切り直しただけです。
その結果、行 1 に、元は別の行にいた 2 と 3 が同居しています。

\begin{figure}[tbp]
\centering
\begin{tikzpicture}[x=5.2mm,y=5.2mm,font=\footnotesize]
  \begin{scope}[shift={(8.5,0.5)}]
    \fill[figblue] (0,1) rectangle (3,2); \fill[figorange] (0,0) rectangle (3,1);
    \draw[black!45,step=1] (0,0) grid (3,2);
    \foreach \v/\x/\y in {0/0/1,1/1/1,2/2/1,3/0/0,4/1/0,5/2/0} \node at (\x+0.5,\y+0.5) {\v};
    \node at (1.5,2.5) {\texttt{A}\ \ 形 $(2,3)$};
  \end{scope}
  \begin{scope}[shift={(2.5,0)}]
    \fill[figblue] (0,0) rectangle (1,3); \fill[figorange] (1,0) rectangle (2,3);
    \draw[black!45,step=1] (0,0) grid (2,3);
    \foreach \v/\x/\y in {0/0/2,3/1/2,1/0/1,4/1/1,2/0/0,5/1/0} \node at (\x+0.5,\y+0.5) {\v};
    \node at (1,3.5) {\texttt{A.T}\ \ 形 $(3,2)$};
    \node[lbl,anchor=east,align=right] at (-0.3,1.5) {行が、そのまま\\列になる};
  \end{scope}
  \begin{scope}[shift={(15.5,0)}]
    \fill[figblue] (0,2) rectangle (2,3); \fill[figblue] (0,1) rectangle (1,2);
    \fill[figorange] (1,1) rectangle (2,2); \fill[figorange] (0,0) rectangle (2,1);
    \draw[black!45,step=1] (0,0) grid (2,3);
    \foreach \v/\x/\y in {0/0/2,1/1/2,2/0/1,3/1/1,4/0/0,5/1/0} \node at (\x+0.5,\y+0.5) {\v};
    \node at (1.6,3.5) {\texttt{A.reshape(3, 2)}\ \ 形 $(3,2)$};
    \node[lbl,anchor=west,align=left] at (2.3,1.5) {2 個ずつ\\区切り直す};
  \end{scope}
  \draw[arr] (8.2,1.5) -- (4.9,1.5) node[midway,above,lbl] {転置};
  \draw[arr] (11.8,1.5) -- (15.1,1.5) node[midway,above,lbl] {reshape};
\end{tikzpicture}
\caption{転置（左）と reshape（右）の違い。色は、元の行列 \texttt{A} での行を表す。
どちらも形は $(3,2)$ になるが、中身の並びは別物である。}
\label{fig:pt-transpose}
\end{figure}

「行と列を入れ替えたい」ときに \code{reshape} を使うのは、初心者に多い間違いです。形は望みどおりになり、エラーも出ませんが、中身は混ざっています。
\textbf{軸を入れ替えるときは転置、区切り方を変えるときは reshape}、と使い分けてください。

ここまでの要点は、形が、1 列に並んだ数の「読み方」だということです。
区切り方を変える \code{reshape} と、軸を入れ替える転置は、形が同じになっても、中身は別物です。

\subsection*{軸を入れ替える・足す・取り除く操作の一覧}
\laterread{初読では飛ばしてかまいません。\chapref{ch:attention}や\chapref{ch:transformer}のコードに \code{view} や \code{transpose} が出てきたら、\tabref{tab:pt-shape-ops}に戻ってください。飛ばしても、次の節は読めます。}

3 階以上のテンソルにも、軸を入れ替える操作や、軸を足したり取り除いたりする操作があります。
\tabref{tab:pt-shape-ops}に、本書で使うものをまとめました。いま覚える必要はありません。あとの章で出てきたら、この表に戻ってください。
このうち \code{X[None]} は、\textbf{長さ 1 の軸}を先頭に足します。要素数も中身も変わりません。
\code{squeeze} はその逆で、長さ 1 の軸を取り除きます。
長さ 1 の軸は無意味に見えますが、次節のブロードキャストで、形をそろえる重要な役を担います。

\begin{table}[tbp]
\centering
\caption{形を変える操作。\texttt{X} は形 $(2,3,4)$ のテンソルとする。\texttt{view} は \texttt{reshape} とほぼ同じ働きで、数を動かさずに区切り方だけを変える。ただし、転置の直後には使えないことがある（\chapref{ch:attention}）。行列の \texttt{A.T} は \texttt{A.transpose(0, 1)} と同じである。}
\label{tab:pt-shape-ops}
\small
\begin{tabular}{@{}lll@{}}
\toprule
コード & 働き & 結果の形 \\
\midrule
\texttt{X.reshape(6, 4)}、\texttt{X.view(6, 4)} & 区切り方を変える & $(6,4)$ \\
\texttt{X.transpose(1, 2)} & 軸 1 と軸 2 を入れ替える & $(2,4,3)$ \\
\texttt{X.permute(2, 0, 1)} & 軸を「2, 0, 1」の順に並べ替える & $(4,2,3)$ \\
\texttt{X.unsqueeze(0)}、\texttt{X[None]} & 先頭に長さ 1 の軸を足す & $(1,2,3,4)$ \\
\texttt{X[None].squeeze(0)} & 長さ 1 の軸 0 を取り除く & $(2,3,4)$ \\
\bottomrule
\end{tabular}
\end{table}

この節では、形を変える二つの操作、\code{reshape} と転置を学びました。残りの操作は、\tabref{tab:pt-shape-ops}から引けます。

\section{まとめて計算する：集計とブロードキャスト}\label{sec:pt-broadcast}

形を読み、変えられるようになりました。いよいよ計算です。
「行ごとに合計する」「各行から、その行の平均を引く」といった計算は、\code{for} 文なしでどう書くのでしょうか。

基本の規則は、同じ形のテンソルどうしの \code{+} や \code{*} が、\textbf{同じ位置の要素どうし}で計算されることです。
$\vx=(1,2,3)$、$\vy=(4,5,6)$ なら、\code{x * y} は $(1\times 4,\,2\times 5,\,3\times 6)=(4,10,18)$ です。内積ではありません（内積 \code{x @ y} なら、この三つを足した 32 です）。
\code{x ** 2}（2 乗）や \code{torch.exp(x)}（指数関数 $e^{x}$）などの関数も、要素ごとに適用されます。

次は、合計や平均です。
例として、2 人の生徒が 3 科目の試験を受けた点数の表（形 $(2,3)$、行が人、列が科目）を使います。

\snippet{軸を指定して集計する}{code:pt-reduce}{ch02/reduce.py}{reduce}

\outputfile{ch02/out/reduce_2.txt}

3 行目の \code{S.sum()} は、全要素の合計で、結果は数 1 個です。$70+80+90+50+60+40=390$ です。
4 行目の \code{S.sum(dim=0)} は 3 個の数（科目ごとの合計）、5 行目の \code{S.sum(dim=1)} は 2 個の数（人ごとの合計）になりました。
\code{dim}（dimension の略で、軸の番号を表します）は、何を指定しているのでしょうか（\figref{fig:pt-dim}）。
\code{dim=0} は、「軸 0 の方向に足し合わせて、\textbf{軸 0 をつぶす}」という指定です。
軸 0 は縦（人の方向）なので、2 人分が 1 個にまとまり、科目ごとの合計が残ります。
科目 0 は $70+50=120$、科目 1 は $80+60=140$、科目 2 は $90+40=130$ です。
形で言えば、$(2,3)$ から軸 0 の「2」が消えて、$(3,)$ になります。
同じように、\code{dim=1} は軸 1（横）をつぶし、$(2,3)$ から「3」が消えて $(2,)$ になります。
人 0 は $70+80+90=240$、人 1 は $50+60+40=150$ です。
「\textbf{指定した軸が、結果の形から消える}」と覚えてください。3 階以上でも、そのまま通用します。
なお、\code{dim=-1} と書くと「最後の軸」を指します。

\begin{figure}[tbp]
\centering
\begin{tikzpicture}[x=8.5mm,y=5.8mm,font=\footnotesize]
  \fill[figblue] (0,0) rectangle (3,2); \draw[black!45,step=1] (0,0) grid (3,2);
  \foreach \v/\x/\y in {70/0/1,80/1/1,90/2/1,50/0/0,60/1/0,40/2/0} \node at (\x+0.5,\y+0.5) {\v};
  \node[lbl,anchor=east] at (-0.1,1.5) {人 0}; \node[lbl,anchor=east] at (-0.1,0.5) {人 1};
  \foreach \c in {0,1,2} \node[lbl,font=\scriptsize] at (\c+0.5,2.4) {科目 \c};
  \node[anchor=east] at (-1.0,1) {\texttt{S}\ \ 形 $(2,3)$};
  \draw[arr,bookaccent] (1.5,-0.2) -- (1.5,-1.2);
  \node[anchor=west,text=bookaccent] at (1.65,-0.7) {軸 0（縦）をつぶす};
  \begin{scope}[shift={(0,-2.4)}]
    \fill[figred] (0,0) rectangle (3,1); \draw[black!45,step=1] (0,0) grid (3,1);
    \foreach \v/\x in {120/0,140/1,130/2} \node at (\x+0.5,0.5) {\v};
    \node[anchor=west] at (3.2,0.5) {\texttt{S.sum(dim=0)}\ \ 形 $(3,)$};
  \end{scope}
  \draw[arr,bookgreen] (3.2,1) -- (5.6,1) node[midway,above,text=bookgreen,align=center] {軸 1（横）\\をつぶす};
  \begin{scope}[shift={(5.8,0)}]
    \fill[figgreen] (0,0) rectangle (1,2); \draw[black!45,step=1] (0,0) grid (1,2);
    \node at (0.5,1.5) {240}; \node at (0.5,0.5) {150};
    \node[anchor=west] at (1.2,1) {\texttt{S.sum(dim=1)}\ \ 形 $(2,)$};
  \end{scope}
\end{tikzpicture}
\caption{\texttt{dim} は「つぶす軸」を指定する。結果の形は、元の形から、指定した軸を取り除いたものになる。}
\label{fig:pt-dim}
\end{figure}

6 行目の \code{S.mean(dim=1)} は、人ごとの平均点です。人ごとの合計を科目の数 3 で割って、$240\div 3=80$ 点と $150\div 3=50$ 点です。
7 行目は、同じ計算に \code{keepdim=True} を付けたものです。数は同じですが、角括弧が二重になっています。
\code{keepdim=True} を付けると、つぶした軸が長さ 1 の軸として残り、形は $(2,)$ ではなく $(2,1)$ になります。
この違いが何の役に立つのかは、すぐあとで分かります。
8 行目の \code{S.argmax(dim=1)} は、軸 1 に沿って、最大値の\textbf{位置}を返します。
結果の \code{[2, 1]} は、「人 0 の最高点は科目 2、人 1 の最高点は科目 1」と読みます。
モデルの予測から「確率が最大のトークンの番号」を取り出すときに使います。
位置ではなく、最大値そのものが欲しいときは、\code{S.max(dim=1).values} と書きます。
\code{max} は、最大値（\code{values}）と、その位置（\code{indices}）の二つを組にして返すからです。この例では $(90,60)$ になります。

\subsection*{ブロードキャスト：形の違う相手と計算する}

集計ができると、「各生徒の点数から、その生徒の平均点を引く」といった計算がしたくなります。
点数の表は形 $(2,3)$、平均点は 2 個の数で、形が違います。
同じ形どうしでしか計算できないなら、平均点を 3 回ずつ並べた $(2,3)$ の表を、わざわざ作ることになります。
PyTorch（と NumPy）は、これを自動でやってくれます。\term{ブロードキャスト}{broadcasting}という仕組みです。
形の違うテンソルどうしの演算で、\textbf{小さいほうを、同じ数の繰り返しで引き伸ばしてから}計算します。

\snippet{ブロードキャストの例}{code:pt-broadcast}{ch02/broadcast.py}{basic}

\outputfile{ch02/out/broadcast_1.txt}

四つの例を、\figref{fig:pt-broadcast}と見比べながら読みます（出力は、2 行ずつが一つの結果です）。
3 行目の \code{S + 5} は、表と数 1 個の足し算です。5 が全部の要素に配られます。
4〜5 行目は、形 $(2,3)$ の表に、形 $(3,)$ の \code{bonus}（科目ごとの加点）を足します。
図の左のとおり、この 1 行が下に向かって複製され、どの生徒にも同じ加点が足されます。
6〜7 行目が、やりたかった計算です。
\code{mean} は \code{keepdim=True} で作ったので、形は $(2,1)$、つまり 2 行 1 列の縦長の表です。
図の中央のとおり、この 1 列が右に向かって複製され、各行から、その行の平均（80 と 50）が引かれます。
出力の 3 番目の表は、行ごとに足すと 0 になっています。
8〜10 行目は、両方が引き伸ばされる例です。
形 $(2,1)$ の縦の列と、形 $(1,3)$ の横の行を足すと、図の右のように、どちらも $(2,3)$ に広げられます。
結果は、すべての組み合わせの和です。

\begin{figure}[tbp]
\centering
\begin{tikzpicture}[x=4.6mm,y=4.6mm,font=\scriptsize,g/.style={text=bookgray}]
  \begin{scope}
    \fill[figblue] (0,0) rectangle (3,2); \draw[black!45,step=1] (0,0) grid (3,2);
    \foreach \v/\x/\y in {70/0/1,80/1/1,90/2/1,50/0/0,60/1/0,40/2/0} \node at (\x+0.5,\y+0.5) {\v};
    \node[font=\small] at (3.5,1) {$+$};
    \fill[figorange] (4,1) rectangle (7,2);
    \draw[black!45,dashed,step=1] (4,0) grid (7,1);
    \draw[black!45,step=1] (4,1) grid (7,2);
    \foreach \v/\x in {0/4,10/5,20/6}{\node at (\x+0.5,1.5) {\v}; \node[g] at (\x+0.5,0.5) {\v};}
    \draw[darr] (7.45,1.5) -- (7.45,0.5);
    \node[font=\footnotesize] at (3.5,-0.9) {$(2,3)+(3,)$};
  \end{scope}
  \begin{scope}[shift={(9.7,0)}]
    \fill[figblue] (0,0) rectangle (3,2); \draw[black!45,step=1] (0,0) grid (3,2);
    \foreach \v/\x/\y in {70/0/1,80/1/1,90/2/1,50/0/0,60/1/0,40/2/0} \node at (\x+0.5,\y+0.5) {\v};
    \node[font=\small] at (3.5,1) {$-$};
    \fill[figorange] (4,0) rectangle (5,2);
    \draw[black!45,dashed,step=1] (5,0) grid (7,2);
    \draw[black!45,step=1] (4,0) grid (5,2);
    \foreach \x in {4,5,6}{\ifnum\x=4 \node at (\x+0.5,1.5) {80}; \node at (\x+0.5,0.5) {50};
      \else \node[g] at (\x+0.5,1.5) {80}; \node[g] at (\x+0.5,0.5) {50}; \fi}
    \draw[darr] (4.5,2.45) -- (6.5,2.45);
    \node[font=\footnotesize] at (3.5,-0.9) {$(2,3)-(2,1)$};
  \end{scope}
  \begin{scope}[shift={(19.4,0)}]
    \fill[figblue] (0,0) rectangle (1,2);
    \draw[black!45,dashed,step=1] (1,0) grid (3,2);
    \draw[black!45,step=1] (0,0) grid (1,2);
    \foreach \x in {0,1,2}{\ifnum\x=0 \node at (\x+0.5,1.5) {1}; \node at (\x+0.5,0.5) {2};
      \else \node[g] at (\x+0.5,1.5) {1}; \node[g] at (\x+0.5,0.5) {2}; \fi}
    \draw[darr] (0.5,2.45) -- (2.5,2.45);
    \node[font=\small] at (3.5,1) {$+$};
    \fill[figorange] (4,1) rectangle (7,2);
    \draw[black!45,dashed,step=1] (4,0) grid (7,1);
    \draw[black!45,step=1] (4,1) grid (7,2);
    \foreach \v/\x in {10/4,20/5,30/6}{\node at (\x+0.5,1.5) {\v}; \node[g] at (\x+0.5,0.5) {\v};}
    \draw[darr] (7.45,1.5) -- (7.45,0.5);
    \node[font=\footnotesize] at (3.5,-0.9) {$(2,1)+(1,3)$};
  \end{scope}
\end{tikzpicture}
\caption{ブロードキャスト。実線のマスが実際に持っている数、破線のマスが、計算のために複製されたとみなされる数である。}
\label{fig:pt-broadcast}
\end{figure}

どんな形の組み合わせなら、引き伸ばしてもらえるのでしょうか。規則は一つだけです。

\begin{definition}[ブロードキャストの規則]
二つのテンソルの形を\textbf{右端をそろえて}並べ、右の軸から順に長さを比べます。
(1) 長さが同じなら、そのまま計算します。
(2) どちらかの長さが 1 なら、その軸を、相手の長さになるまで複製します。
(3) 片方に軸がなければ、そこに長さ 1 の軸があるものとみなします。
(4) 長さが異なり、どちらも 1 でなければ、エラーになります。
結果の形は、各軸で長いほうを取ったものです。
\end{definition}

\figref{fig:pt-bc-rule}は、この規則を四つの例に当てはめたものです。
1 番目は、先ほどの \code{S + bonus} です。右端の「3 と 3」は同じなので、そのままです。
その左は、$(2,3)$ の「2」に対して、\code{bonus} には軸がありません。長さ 1 とみなされ、2 に引き伸ばされます。
形 $(3,)$ の \code{bonus} が「横の 1 行」としてふるまったのは、右端でそろえるからです。
2 番目は \code{col + row} で、どちらの軸でも、長さ 1 の側が引き伸ばされます。

\begin{figure}[tbp]
\centering
\begin{tikzpicture}[x=6.5mm,y=5.4mm,font=\footnotesize,
  d/.style={draw=black!40,minimum width=5.6mm,minimum height=4.5mm,inner sep=0pt},
  ok/.style={d,fill=figgreen},st/.style={d,fill=figorange},ng/.style={d,fill=figred},
  no/.style={d,dashed,fill=figorange,text=bookgray}]
  \begin{scope}
    \node[anchor=east] at (0.5,2) {\texttt{S}}; \node[d] at (1,2) {2}; \node[ok] at (2,2) {3};
    \node[anchor=east] at (0.5,1) {\texttt{bonus}}; \node[no] at (1,1) {1}; \node[ok] at (2,1) {3};
    \draw[black!60] (0.55,0.5) -- (2.45,0.5);
    \node[anchor=east] at (0.5,0) {結果}; \node[d] at (1,0) {2}; \node[d] at (2,0) {3};
  \end{scope}
  \begin{scope}[shift={(4.6,0)}]
    \node[anchor=east] at (0.5,2) {\texttt{col}}; \node[d] at (1,2) {2}; \node[st] at (2,2) {1};
    \node[anchor=east] at (0.5,1) {\texttt{row}}; \node[st] at (1,1) {1}; \node[d] at (2,1) {3};
    \draw[black!60] (0.55,0.5) -- (2.45,0.5);
    \node[anchor=east] at (0.5,0) {結果}; \node[d] at (1,0) {2}; \node[d] at (2,0) {3};
  \end{scope}
  \begin{scope}[shift={(9.2,0)}]
    \node[anchor=east] at (0.5,2) {\texttt{S}}; \node[d] at (1,2) {2}; \node[ng] at (2,2) {3};
    \node[anchor=east] at (0.5,1) {\texttt{m}}; \node[no] at (1,1) {1}; \node[ng] at (2,1) {2};
    \draw[black!60] (0.55,0.5) -- (2.45,0.5);
    \node[anchor=east] at (0.5,0) {結果}; \node[text=bookaccent] at (1.5,0) {エラー};
  \end{scope}
  \begin{scope}[shift={(13.8,0)}]
    \node[anchor=east] at (0.5,2) {\texttt{pred}}; \node[d] at (1,2) {3}; \node[st] at (2,2) {1};
    \node[anchor=east] at (0.5,1) {\texttt{target}}; \node[no] at (1,1) {1}; \node[d] at (2,1) {3};
    \draw[black!60] (0.55,0.5) -- (2.45,0.5);
    \node[anchor=east] at (0.5,0) {結果}; \node[d] at (1,0) {?}; \node[d] at (2,0) {?};
  \end{scope}
\end{tikzpicture}
\caption{ブロードキャストの規則の当てはめ方。形を右端でそろえて縦に並べ、右の列から順に比べる。
緑は同じ長さ、橙は引き伸ばされる長さ 1 の軸、破線は軸がなく 1 とみなされる場所、赤は長さが合わない組である。4 番目の結果は、本文で考える。}
\label{fig:pt-bc-rule}
\end{figure}

3 番目は、失敗する例です。
先ほどの「平均を引く」計算で \code{keepdim=True} を付け忘れると、平均 \code{m} の形は $(2,)$ です。
右端をそろえると、$(2,3)$ の「3」と $(2,)$ の「2」が並びます。
長さが違い、どちらも 1 ではないので、エラーになります（エラー文の読み方は、\secref{sec:pt-final}で扱います）。
「2 人分の平均が 2 個あるのだから、行ごとに引いてほしい」と思っても、規則は右端からしか比べません。
\code{keepdim=True} で形を $(2,1)$ にしておけば、右端は「3 と 1」になり、正しく引き伸ばされます。
これが、\code{keepdim} の役目です。

\subsection*{黙って起こる間違い}

エラーが出るのは、じつは幸運な場合です。
困るのは、\textbf{形がたまたま合ってしまい、間違った計算が黙って進む}場合です。
\figref{fig:pt-bc-rule}の 4 番目が、その例です。
モデルの予測は形 $(N,1)$、正解は形 $(N,)$ で手に入ることが多く、この組み合わせは実際によく起こります。
データ 1 個につき数を 1 個だけ予測するモデルでも、$N$ 個分の出力は「$N$ 行 1 列の行列」として返ってくるからです（理由は\secref{sec:pt-matmul}で分かります）。

\begin{quiz}
形 $(3,1)$ のテンソル \code{pred} から、形 $(3,)$ のテンソル \code{target} を引きます。
結果の形は何になるでしょうか。規則を当てはめて考えてください。
\end{quiz}

右端は「1 と 3」です。片方が 1 なので、3 に引き伸ばされます。
その左は「3 と、軸なし」です。軸なしは 1 とみなされ、3 に引き伸ばされます。
結果の形は $(3,3)$ です。
3 個の予測と 3 個の正解の差を取るつもりなら、欲しいのは 3 個の数のはずです。実行して確かめます。

\snippet{エラーにならない間違い}{code:pt-bug-outer}{ch02/broadcast_bug.py}{outer}

\outputfile{ch02/out/broadcast_bug_2.txt}

\code{pred} と \code{target} には、同じ値 $(1,2,3)$ を入れました。予測が完全に当たっているので、差は全部 0 のはずです。
ところが、差 \code{diff} の形は、出力の 1 行目のとおり $(3,3)$ です。
中身は\figref{fig:pt-bug}のとおりで、予測の一つひとつが、\textbf{すべての}正解と引き算されています。
欲しかったのは、対角線上の 3 個だけです。
5 行目で 2 乗して平均すると、0 になるはずの誤差が 1.3333 になりました。
図の 9 個の差を 2 乗すると、0 が 3 個、1 が 4 個、4 が 2 個です。合計は $1\times 4+4\times 2=12$ で、9 で割ると $1.3333$ です。
エラーは出ません。学習に使えば、間違った目標に向かって、黙って学習が進みます。
直し方は、6 行目です。\code{pred.squeeze(1)} は、長さ 1 の軸 1 を取り除く操作で、形 $(3,1)$ を $(3,)$ にします。
こうして形を $(3,)$ にそろえてから引くと、誤差は正しく 0 になります。

\begin{figure}[tbp]
\centering
\begin{tikzpicture}[x=6mm,y=6mm,font=\footnotesize]
  \begin{scope}[shift={(0,3.4)}]
    \fill[figorange] (0,0) rectangle (3,1); \draw[black!45,step=1] (0,0) grid (3,1);
    \foreach \v/\x in {1/0,2/1,3/2} \node at (\x+0.5,0.5) {\v};
  \end{scope}
  \node[anchor=west] at (3.2,3.9) {\texttt{target}\ \ 形 $(3,)$};
  \begin{scope}[shift={(-1.4,0)}]
    \fill[figblue] (0,0) rectangle (1,3); \draw[black!45,step=1] (0,0) grid (1,3);
    \foreach \v/\y in {1/2,2/1,3/0} \node at (0.5,\y+0.5) {\v};
  \end{scope}
  \node[anchor=east] at (-1.6,1.5) {\texttt{pred}\ \ 形 $(3,1)$};
  \fill[figred] (0,0) rectangle (3,3);
  \fill[figgreen] (0,2) rectangle (1,3); \fill[figgreen] (1,1) rectangle (2,2); \fill[figgreen] (2,0) rectangle (3,1);
  \draw[black!45,step=1] (0,0) grid (3,3);
  \foreach \v/\x/\y in {0/0/2,-1/1/2,-2/2/2,1/0/1,0/1/1,-1/2/1,2/0/0,1/1/0,0/2/0} \node at (\x+0.5,\y+0.5) {$\v$};
  \node[anchor=west,align=left] at (3.3,1.5) {\texttt{pred - target}\ \ 形 $(3,3)$\\[0.5mm]
    緑: 欲しかった 3 個（対応する相手との差）\\赤: 無関係な相手との差};
\end{tikzpicture}
\caption{形 $(3,1)$ から形 $(3,)$ を引くと、「すべての組み合わせの差」が計算され、形 $(3,3)$ になる。エラーにならないので、気づきにくい。}
\label{fig:pt-bug}
\end{figure}

こうした間違いを防ぐ習慣が、二つあります。
一つは、集計の結果を元のテンソルとの計算に使うときに、\code{keepdim=True} を付けることです。
もう一つは、コードの行末に、形をコメントで書くことです（\texttt{\# (2, 3)} のように）。本書のコードも、そうしています。

この節では、\code{dim} で軸を指定する集計と、ブロードキャストを学びました。便利な反面、形の取り違えがエラーにならないことがあります。

\subsection*{ここまでの道のり}

\begin{itemize}
\item 要素を 1 個ずつ回る \code{for} 文は、配列演算の 1,000 倍以上遅い。だから、テンソルにまとめて計算する。
\item テンソルは、1 列に並んだ数と、その読み方である「形」、数の持ち方である「型」からできている。
\item 形は \code{reshape} や転置で変えられる。形の違う相手との計算は、ブロードキャストが引き伸ばしてくれる。
\end{itemize}

残る道具は二つです。モデルの計算の主役である\textbf{行列積}と、学習の主役である\textbf{自動微分}です。

\section{行列積と「1 データ 1 行」の約束}\label{sec:pt-matmul}

前節の演算は、どれも「同じ位置の要素どうし」の計算でした。
しかし、モデルの中で最も多く使われるのは、そうではない計算、つまり行列積です。
\chapref{ch:math}では、$\vy=\mW\vx+\vb$ という式で、ベクトルを別のベクトルに移しました。
この式を、128 個のトークンを表す 128 個のベクトルのような、たくさんのデータに一度に適用するには、どう書けばよいのでしょうか。

まず、書き方の確認です。行列積は、コードでは \code{A @ B} と書きます。
形の規則は、\chapref{ch:math}で学んだとおりです。
\textbf{左の形の右端と、右の形の左端が一致}していれば計算でき、結果の形は、残った外側の二つです。
$(3,\underline{2})$ と $(\underline{2},4)$ なら、内側の 2 どうしが一致し、結果は $(3,4)$ です。
相手はベクトルでもかまいません。形 $(3,\underline{2})$ の \code{A} と形 $(\underline{2},)$ の \code{v} なら、内側の 2 どうしが一致します。
残る外側は 3 だけなので、\code{A @ v} は形 $(3,)$ のベクトルです。数式の $\mA\vv$ にあたります。
一方、\code{A * B} は要素ごとの積です。\textbf{行列積は \code{@}} です（\file{ch02/matmul.py}）。
なお、積を添字で書く \code{torch.einsum} という関数もあり、同じファイルで試せます。

ここで、本書全体に関わる約束を一つ説明します。
数式では、ベクトルを縦に並んだ\textbf{列ベクトル}として扱い、行列を左から掛けて $\vy=\mW\vx+\vb$ と書きます。
3 個の数からなる $\vx$ を、2 個の数からなる $\vy$ に移すなら、$\mW$ は 2 行 3 列です。
$\mW\vx$ の成分は、$\mW$ の行と $\vx$ の内積を、行ごとに並べたものでした。
だから、$\mW$ の列の数は $\vx$ の長さ 3 に合わせ、行の数は $\vy$ の長さ 2 に合わせます。
一方、コードでは、$N$ 個のデータ $\vx^{(1)},\dots,\vx^{(N)}$ を、まとめて 1 個のテンソルに入れます。
このとき、\textbf{1 個のデータを 1 行として、縦に積む}のが、PyTorch の決まりです。
データが 4 個、それぞれが 3 個の数なら、形 $(4,3)$ の行列 $\mX$ になります。軸 0 が「何個目のデータか」を表します。
こうしてまとめたデータの束を、\term{バッチ}{batch}と呼びます。

欲しいのは、$\mX$ の各行に $\mW$ を掛けた結果を、同じく行として積んだ、形 $(4,2)$ の行列 $\mY$ です。
どんな行列積で得られるでしょうか。形から考えます。
$\mX$ は $(4,3)$、$\mW$ は $(2,3)$ です。
\code{W @ X} は、内側が「3 と 4」で合いません。\code{X @ W} も、内側が「3 と 2」で合いません。
$\mW$ を転置して $(3,2)$ にすると、$(4,\underline{3})$ と $(\underline{3},2)$ で内側が合い、結果は望みどおり $(4,2)$ になります。
中身も正しいことを、1 マスだけ確かめます。
$\mX$ の先頭の行（1 個目のデータ）を $(1,2,3)$、$\mW$ の先頭の行を $(1,0,-1)$ とします。
転置すると、行は列になります。$\mW$ の先頭の行は、$\mW^\T$ では先頭の列です。
だから、$\mX\mW^\T$ の左上のマスは、$(1,2,3)$ と $(1,0,-1)$ の内積 $1+0-3=-2$ です。
これは、数式の $\mW\vx$ の最初の成分（$\mW$ の先頭の行と $\vx$ の内積）と、同じ計算です。
ほかのマスも同じで、$\mX\mW^\T$ の各行は、その行のデータに $\mW$ を掛けた結果になります。
まとめると、次の対応になります（\figref{fig:pt-rows}）。
\begin{equation}
\underbrace{\vy=\mW\vx+\vb}_{\text{数式：1 個のデータ（列ベクトル）}}
\qquad\Longleftrightarrow\qquad
\underbrace{\texttt{Y = X @ W.T + b}}_{\text{コード：$N$ 個のデータ（1 データ 1 行）}}
\label{eq:pt-rows}
\end{equation}
最後の \code{+ b} は、ブロードキャストです。
形 $(4,2)$ に形 $(2,)$ を足すので、どの行にも同じ $\vb$ が足されます。
なお、出力が数 1 個のモデルでは、$\mW$ は 1 行 3 列です。$(4,\underline{3})$ と $(\underline{3},1)$ の積なので、$\mY$ の形は $(4,)$ ではなく $(4,1)$ になります。
\secref{sec:pt-broadcast}で「予測は形 $(N,1)$ で手に入ることが多い」と述べたのは、このためです。

\begin{figure}[tbp]
\centering
\begin{tikzpicture}[x=4.3mm,y=4.3mm,font=\footnotesize]
  \node[anchor=west] at (-0.5,6.3) {\textbf{数式}：データ 1 個は列ベクトル};
  \begin{scope}[shift={(0,1.5)}]
    \fill[figgreen] (0,0.5) rectangle (1,2.5);
    \draw[black!45] (0,0.5) rectangle (1,2.5); \draw[black!45] (0,1.5) -- (1,1.5);
    \node at (0.5,3.1) {$\vy$};
    \node at (1.8,1.5) {$=$};
    \fill[figorange] (2.6,0.5) rectangle (5.6,2.5); \fill[bookorange!55] (2.6,1.5) rectangle (5.6,2.5);
    \draw[black!45] (2.6,0.5) rectangle (5.6,2.5); \draw[black!45] (2.6,1.5) -- (5.6,1.5);
    \draw[black!45] (3.6,0.5) -- (3.6,2.5); \draw[black!45] (4.6,0.5) -- (4.6,2.5);
    \node at (4.1,3.1) {$\mW$\ $(2,3)$};
    \fill[figblue] (6.3,0) rectangle (7.3,3);
    \draw[black!45] (6.3,0) rectangle (7.3,3); \draw[black!45] (6.3,1) -- (7.3,1); \draw[black!45] (6.3,2) -- (7.3,2);
    \node at (6.8,3.6) {$\vx$};
    \node at (8.1,1.5) {$+$};
    \fill[figpurple] (8.9,0.5) rectangle (9.9,2.5);
    \draw[black!45] (8.9,0.5) rectangle (9.9,2.5); \draw[black!45] (8.9,1.5) -- (9.9,1.5);
    \node at (9.4,3.1) {$\vb$};
  \end{scope}
  \begin{scope}[shift={(14,0)}]
    \node[anchor=west] at (-0.5,6.3) {\textbf{コード}：データ 1 個は 1 行};
    \fill[figgreen] (0,0.5) rectangle (2,4.5); \fill[bookgreen!50] (0,2.5) rectangle (2,3.5);
    \draw[black!45] (0,0.5) rectangle (2,4.5); \draw[black!45] (1,0.5) -- (1,4.5);
    \foreach \y in {1.5,2.5,3.5} \draw[black!45] (0,\y) -- (2,\y);
    \node at (1,5.1) {\texttt{Y}\ $(4,2)$};
    \node at (2.8,2.5) {$=$};
    \fill[figblue] (3.6,0.5) rectangle (6.6,4.5); \fill[bookblue!45] (3.6,2.5) rectangle (6.6,3.5);
    \draw[black!45] (3.6,0.5) rectangle (6.6,4.5); \draw[black!45] (4.6,0.5) -- (4.6,4.5); \draw[black!45] (5.6,0.5) -- (5.6,4.5);
    \foreach \y in {1.5,2.5,3.5} \draw[black!45] (3.6,\y) -- (6.6,\y);
    \node at (5.1,5.1) {\texttt{X}\ $(4,3)$};
    \node at (7.3,2.5) {\texttt{@}};
    \fill[figorange] (8,1) rectangle (10,4); \fill[bookorange!55] (8,1) rectangle (9,4);
    \draw[black!45] (8,1) rectangle (10,4); \draw[black!45] (9,1) -- (9,4); \draw[black!45] (8,2) -- (10,2); \draw[black!45] (8,3) -- (10,3);
    \node at (9,4.6) {\texttt{W.T}\ $(3,2)$};
    \node at (10.8,2.5) {$+$};
    \fill[figpurple] (11.6,2) rectangle (13.6,3);
    \draw[black!45] (11.6,2) rectangle (13.6,3); \draw[black!45] (12.6,2) -- (12.6,3);
    \node at (12.6,3.6) {\texttt{b}\ $(2,)$};
    \node[lbl,anchor=west,align=left] at (-0.5,-0.8) {濃い行: 1 個のデータ $\vx^{(i)}$ と、その結果 $\vy^{(i)}$\\濃い橙: $\mW$ の 1 行は、\texttt{W.T} では 1 列になる};
  \end{scope}
\end{tikzpicture}
\caption{数式の書き方（左）と、コードの書き方（右）。数式ではベクトルを縦に置き、行列を左から掛ける。
コードでは、データを 1 個 1 行として縦に積み、転置した行列を右から掛ける。}
\label{fig:pt-rows}
\end{figure}

両方の書き方で計算して、一致することを確かめます。

\snippet{1 個ずつの計算と、まとめての計算}{code:pt-rows}{ch02/matmul.py}{rows}

\outputfile{ch02/out/matmul_2.txt}

1〜3 行目で $\mW$ と $\vb$ を、4〜7 行目で 4 個のデータを積んだ $\mX$ を作ります。
8〜9 行目は、数式どおりの計算です。
\code{for x in X} で \code{X} の行を 1 個ずつ取り出し、それぞれに \code{W @ x + b} を計算します。
出力の最初の 4 行が、その結果です。
1 個目は、$\vx=(1,2,3)$ に対して $\mW\vx=(1-3,\ 2+2)=(-2,4)$、これに $\vb=(0.5,-0.5)$ を足して $(-1.5,3.5)$ です。
10 行目が、式\eqref{eq:pt-rows}の右側です。\code{for} 文なしの 1 行で、4 個まとめて計算します。
出力の後半の行列は、行ごとに見ると、前半の 4 行と完全に一致しています。

この約束は、PyTorch 自身も使っています。
\secref{sec:pt-tensor}のモデルの一覧にあった形 $(576,192)$ のテンソルは、192 個の数を 576 個の数に移す行列 $\mW$ です。
PyTorch は、$\mW$ を数式と同じ「（出力の個数, 入力の個数）」の形で持ち、計算のときに \code{x @ W.T} を実行します。

実際のモデルでは、軸がさらに増えます。
文が 8 個、各文にトークンが 5 個、各トークンが 3 個の数で表されていれば、形は $(8,5,3)$ です。
3 階以上のテンソルに \code{@} を使うと、\textbf{最後の二つの軸}を行列とみなし、それより左の軸は「行列が何枚あるか」として扱います。
形 $(8,5,3)$ の \code{H} に \code{H @ W.T + b} を計算すると、結果の形は $(8,5,2)$ です（\file{ch02/matmul.py}）。
\figref{fig:pt-batch-matmul}が、そのようすです。
5 行 3 列の行列が 8 枚あり、その 1 枚 1 枚に、同じ \code{W.T} が掛かります。
1 枚分は、これまでと同じ規則で、$(5,\underline{3})$ と $(\underline{3},2)$ から $(5,2)$ です。左の軸の 8 は、結果にそのまま残ります。
最後の \code{+ b} は、形 $(8,5,2)$ と形 $(2,)$ のブロードキャストです。
結局、$8\times 5=40$ 個のベクトル全部に、同じ $\mW$ と $\vb$ が適用されます。式\eqref{eq:pt-rows}の書き方は、軸が増えてもそのまま使えます。

\begin{figure}[tbp]
\centering
\begin{tikzpicture}[x=4.3mm,y=4.3mm,font=\footnotesize]
  \foreach \k in {2,1,0}{
    \begin{scope}[shift={(0.55*\k,0.55*\k)}]
      \fill[figblue] (0,0) rectangle (3,5);
      \draw[black!45,step=1] (0,0) grid (3,5);
    \end{scope}}
  \node at (2,6.9) {\texttt{H}\ \ $(8,5,3)$};
  \node[lbl,anchor=east] at (-0.2,2.5) {5 行};
  \node[lbl] at (1.5,-0.6) {3 列};
  \draw[arr,bookgreen] (3.25,-0.1) -- (4.35,1.0);
  \node[lbl,anchor=west,text=bookgreen] at (3.9,-0.1) {8 枚};
  \node at (5.0,3) {\texttt{@}};
  \begin{scope}[shift={(7.4,1.5)}]
    \fill[figorange] (0,0) rectangle (2,3);
    \draw[black!45,step=1] (0,0) grid (2,3);
    \node at (1,3.7) {\texttt{W.T}\ \ $(3,2)$};
    \node[lbl,anchor=east] at (-0.1,1.5) {3 行};
    \node[lbl] at (1,-0.6) {2 列};
  \end{scope}
  \node at (10.6,3) {$=$};
  \begin{scope}[shift={(12,0)}]
    \foreach \k in {2,1,0}{
      \begin{scope}[shift={(0.55*\k,0.55*\k)}]
        \fill[figgreen] (0,0) rectangle (2,5);
        \draw[black!45,step=1] (0,0) grid (2,5);
      \end{scope}}
    \node at (1.6,6.9) {結果\ \ $(8,5,2)$};
    \node[lbl] at (1,-0.6) {2 列};
  \end{scope}
  \node[anchor=west,align=left] at (16.6,3.2) {形の読み方\\[0.5mm]
    $(8,\ \underline{5,\ 3})$\ \ \texttt{@}\ \ $(\underline{3,\ 2})$\ $\to$\ $(8,\ \underline{5,\ 2})$\\[0.5mm]
    下線の部分: 1 枚ごとの行列積\\左端の 8: 行列が何枚あるか};
\end{tikzpicture}
\caption{3 階のテンソルと行列積。\texttt{@} は、最後の二つの軸（5 行 3 列）を 1 枚の行列とみなし、8 枚のそれぞれに同じ \texttt{W.T} を掛ける。
左端の軸の長さ 8 は、結果にそのまま残る。図では、8 枚のうち 3 枚だけを描いた。}
\label{fig:pt-batch-matmul}
\end{figure}

この節では、行列積の形の規則と、「数式は列ベクトル、コードは 1 データ 1 行」という約束を学びました。

\section{自動微分のさわり}\label{sec:pt-autograd}

モデルの計算に必要な道具は、前節まででそろいました。残るのは、学習のための道具です。
\secref{sec:pt-tensor}で、学習とは「テンソルの中の数を少しずつ書き換える作業」だと述べました。
では、どの数を、どちら向きに書き換えればよいのでしょうか。
それを教えてくれるのが、各変数での偏微分を並べた勾配でした（\chapref{ch:math}）。問題は、その求め方です。
\chapref{ch:math}の数値微分は、変数を 1 個だけ少し動かして、関数の値がどれだけ変わるかを見る方法でした。
変数 1 個につき関数を 2 回計算するので、変数が 260 万個あれば 520 万回です。とても使えません。
かといって、260 万個の偏微分の式を手で導くこともできません。
PyTorch は、この微分を\textbf{自動で}求めてくれます。この節では、その使い方だけを学びます。

小さな例で試します。
関数 $f(x_1,x_2)=x_1^{2}+3x_1x_2$ の、点 $(x_1,x_2)=(2,1)$ での勾配を求めます。
先に手で計算します。$x_1$ で偏微分すると $\pd{f}{x_1}=2x_1+3x_2$ で、値を入れると $2\times 2+3\times 1=7$ です。
$x_2$ で偏微分すると $\pd{f}{x_2}=3x_1$ で、$3\times 2=6$ です。勾配は $(7,6)$ になるはずです。

\snippet{自動微分で勾配を求める}{code:pt-autograd}{ch02/autograd.py}{basic}

\outputfile{ch02/out/autograd_1.txt}

手順は三つです（\figref{fig:pt-autograd}）。

\textbf{(1) 印を付ける}。
1 行目で、テンソルを作るときに \code{requires_grad=True} を付けます。「あとで、この変数で微分する」という印です。

\textbf{(2) 普通に計算する}。
2 行目で、$f$ の式をそのまま書きます。\code{x[0]} が $x_1$、\code{x[1]} が $x_2$ です。
出力の 1 行目では、値 10（$2^{2}+3\times 2\times 1$）の後ろに \code{grad_fn=<AddBackward0>} が付いています。
「この値は足し算の結果として作られた」という記録です。
PyTorch は、印の付いたテンソルを使った計算を、裏ですべて記録しています。

\textbf{(3) \code{backward()} を呼ぶ}。
4 行目の \code{f.backward()} で、PyTorch は記録を逆向きにたどり、$f$ を \code{x} の各成分で偏微分します。
結果は \code{x.grad} に入ります。出力の 2 行目は \code{tensor([7., 6.])} で、手計算の $(7,6)$ と一致しました。

\begin{figure}[tbp]
\centering
\begin{tikzpicture}[x=1mm,y=1mm,font=\footnotesize]
  \node[box,minimum width=24mm,align=center] (x) at (0,0) {\texttt{x}\\$(2,\ 1)$};
  \node[plainbox,minimum width=30mm,minimum height=9mm,align=center] (op) at (44,0) {$x_1^{2}+3x_1x_2$};
  \node[greenbox,minimum width=20mm,align=center] (f) at (86,0) {\texttt{f}\\$10$};
  \draw[arr] (x) -- (op); \draw[arr] (op) -- (f);
  \node[lbl] at (0,9) {(1) 印を付ける};
  \node[lbl,text=bookblue] at (58,9) {(2) 普通に計算する（手順が記録される）};
  \node[redbox,minimum width=24mm,align=center] (g) at (0,-16) {\texttt{x.grad}\\$(7,\ 6)$};
  \draw[darr,draw=bookaccent] (f.south) |- (g.east);
  \node[lbl,text=bookaccent] at (48,-12.5) {(3) \texttt{f.backward()}: 記録を逆向きにたどる};
  \node[lbl,anchor=west] at (14,-19.5) {\texttt{x} と同じ形};
\end{tikzpicture}
\caption{自動微分の使い方。印を付けた変数から値を計算すると、手順が記録される。
\texttt{backward()} は、その記録を逆向きにたどって勾配を求める。}
\label{fig:pt-autograd}
\end{figure}

微分の式は、一度も書いていません。
計算の手順を記録し、そこから微分を機械的に求めるこの仕組みを、\term{自動微分}{automatic differentiation}と呼びます（PyTorch での名前は autograd です）。
\chapref{ch:math}の数値微分で検算すると、$7.00000000002$ と $6.000000000004$ で、自動微分の答えと一致しました（\file{ch02/autograd.py}）。
数値微分は、変数 2 個のために $f$ を 4 回計算しました。
自動微分は、\code{backward()} の 1 回で、全部の変数の偏微分を求めます。変数が 260 万個になっても 1 回で、しかも近似ではありません。
なぜそんなことができるのかは、\chapref{ch:nn}で、自分の手で作りながら解き明かします。

自動微分には、知らないと必ずつまずく決まりが三つあります。

一つ目は、\textbf{勾配は上書きではなく、足し込まれる}ことです。
先ほどの続きで、別の関数 $g=x_1^{2}+x_2^{2}$ の勾配を求めてみます。正しい答えは $(2x_1,2x_2)=(4,2)$ です。
\snippet{勾配は足し込まれる}{code:pt-accumulate}{ch02/autograd.py}{accumulate}

\outputfile{ch02/out/autograd_3.txt}

1 行目の \code{(x ** 2).sum()} は、要素ごとに 2 乗して合計する計算で、$g$ そのものです。
ところが、2 行目の \code{g.backward()} のあとの \code{x.grad} は、出力の 1 行目のとおり $(11,8)$ になります。
\code{x.grad} に残っていた前回の $(7,6)$ に、今回の $(4,2)$ が足されたのです。$7+4=11$、$6+2=8$ です。
4 行目の \code{x.grad.zero_()} で勾配を 0 に戻し、5〜6 行目で同じ計算をやり直すと、出力の 2 行目のとおり、正しく $(4,2)$ になります。
\code{zero_} の末尾の下線は、「新しいテンソルを作らず、その場で中身を書き換える」操作の印です。
足し込む理由は、\chapref{ch:nn}で分かります。
学習では \code{backward()} を何千回も呼ぶので、勾配を 0 に戻す 1 行が、\chapref{ch:ml}以降の学習のコードに必ず現れます。

二つ目は、\textbf{記録が不要な計算は、\code{with torch.no_grad():} の下に、字下げして書く}ことです。
字下げした範囲では、PyTorch は計算の記録を取りません。その分、メモリを使わず、速く動きます。
学習済みのモデルで予測だけをするときに使います。記録がないので、その結果に \code{backward()} は呼べません。

三つ目は、\textbf{出発点は数 1 個、勾配の形は変数と同じ}という決まりです。
\code{backward()} を呼ぶ相手は、数 1 個のテンソルに限ります。変数はどんな形でもよく、形 $(2,3)$ の行列 \code{W} なら、勾配 \code{W.grad} も形 $(2,3)$ です。
\code{W.grad[i, j]} は、「\code{W[i, j]} を少し増やすと、値がどれだけ増えるか」を表します。

この節では、印を付け、普通に計算し、\code{backward()} を呼ぶという、自動微分の使い方を学びました。

\section{仕上げ：エラーの読み方と、本物のモデルでの腕試し}\label{sec:pt-final}

道具はすべて出そろいました。最後に、二つの仕上げをします。
一つは、腕試しです。この章の道具だけで書いた softmax と交差エントロピーは、本物の基盤モデルに通用するでしょうか。
その前に、もう一つの仕上げとして、形のエラーに出会ったときの直し方を見ておきます。

\subsection*{形のエラーに出会ったら}
\laterread{自分でコードを動かして、形のエラーが出たときに戻ってくれば足ります。あとの腕試しに出てくる \code{assert} は、「条件が成り立たなければ、その場で止まる」文だと知っていれば読めます。}

\tabref{tab:pt-errors}は、よくある間違いをわざと起こし、出てきたエラー文を集めたものです（\file{ch02/shape_errors.py}）。
英語の長い文に見えますが、読むべき場所は決まっています。
\textbf{数字を探し、自分のテンソルの形と見比べる}のです。
たとえば 1 番目の \code{(4x3 and 2x3)} は、形 $(4,3)$ と $(2,3)$ の行列積をしようとして、内側の 3 と 2 が合わなかった、という意味です。\secref{sec:pt-matmul}の \code{X @ W.T} で、転置を付け忘れるとこうなります。
エラーが出たら、まず関係するテンソルの \code{.shape} を \code{print} で表示してください。頭の中の形と、どこかで食い違っているはずです。

\begin{table}[tbp]
\centering
\caption{よくあるエラー文の読み方。エラー文の中の数字が、手がかりになる。}
\label{tab:pt-errors}
\footnotesize
\begin{tabularx}{\linewidth}{@{}>{\raggedright\arraybackslash}p{43mm}>{\raggedright\arraybackslash}X>{\raggedright\arraybackslash}p{31mm}@{}}
\toprule
エラー文の要点 & 意味 & 最初に確かめること \\
\midrule
\texttt{shapes cannot be multiplied (4x3 and 2x3)} & 行列積の内側の長さ（3 と 2）が合わない & 転置の付け忘れ、掛ける順番 \\
\texttt{size of tensor a (3) must match} \dots\ \texttt{b (2) at non-singleton dimension 1} & ブロードキャストできない。軸 1 の長さが 3 と 2 で、どちらも 1 でない（non-singleton） & \texttt{keepdim} の付け忘れ \\
\texttt{shape '[5, -1]' is invalid for input of size 12} & 要素数 12 は、5 行には区切れない & 元の形と要素数 \\
\texttt{expected} \dots\ \texttt{same dtype} \dots\ \texttt{float != long int} & 小数（\texttt{float32}）と整数（\texttt{int64}、エラー文では long int）を混ぜた & \texttt{.dtype} を調べる \\
\bottomrule
\end{tabularx}
\end{table}

\secref{sec:pt-broadcast}で見たとおり、本当に怖いのは、エラーにならない間違いでした。
これを防ぐ道具が、Python の \code{assert} 文です。
\code{assert 条件, メッセージ} と書くと、条件が成り立たないときに、メッセージを表示してその場でプログラムが止まります。
関数の入口に \code{assert y.shape == (N,)} のような形の約束を書いておけば、形を取り違えた瞬間に気づけます。

\subsection*{腕試し：softmax と交差エントロピーを書く}
\skimread{結論は「この章の道具で書いた数行の関数が、本物のモデルの損失 3.416 を、表示された全桁まで再現する」です。初読では、\coderef{code:pt-softmax}と\figref{fig:pt-ce-flow}を見て、最後の出力を確かめれば足ります。\coderef{code:pt-ce-real}を 1 行ずつ追うのは、\chapref{ch:lm}を読むときでかまいません。}

\chapref{ch:math}で、ロジット（確率に直す前のスコア）を確率に直す softmax と、予測の悪さを測る交差エントロピーを学びました。
そこでは、1 個の予測を扱いました。ここでは、$N$ 個の予測をまとめて計算します。
それぞれの予測が $V$ 個の候補にロジットを付けているので、ロジットをまとめたテンソル $z$ の形は $(N,V)$ です。
$n$ 番目の予測の正解の番号を $y_n$ とすると、計算したいのは次の二つです。
\begin{equation}
p_{nv}=\frac{e^{z_{nv}}}{\sum_{v'=1}^{V}e^{z_{nv'}}},
\qquad
\calL=-\frac{1}{N}\sum_{n=1}^{N}\log p_{n,y_n}
\label{eq:pt-ce}
\end{equation}
左が softmax で、各行のロジットを、0 以上で、足すと 1 になる数に直します。
右が損失で、正解の候補に割り当てた確率 $p_{n,y_n}$ の対数を取り、符号を変えて平均します。
対数を取って符号を変えるのは、\chapref{ch:math}で学んだ、予測の悪さの測り方です。
正解に確率 1 を付ければ $-\log 1=0$ で、付けた確率が小さいほど、値は大きくなります。

\snippet{softmax と交差エントロピー}{code:pt-softmax}{ch02/softmax_ce.py}{functions}

\code{softmax} の 3 行目で、全要素の指数関数を取ります。
4 行目の \code{e.sum(dim=-1, keepdim=True)} は、最後の軸（候補の軸）をつぶして、行ごとの合計を求めます。
\code{keepdim=True} なので、形は $(N,1)$ です。
形 $(N,V)$ を形 $(N,1)$ で割るので、ブロードキャストにより、各行がその行の合計で割られます。
なお、この素朴な書き方には弱点があります。
ロジットが大きいと、$e^{z}$ が小数で表せる範囲を超えて、無限大（\code{inf}）になります。
無限大を無限大で割ると、結果は \code{nan}（数ではない、という印）です（\chapref{ch:math}）。
この弱点は、章末の演習で直します。

\code{cross_entropy} では、まず 9 行目で、\code{z.shape} の二つの数を \code{N} と \code{V} に受け取ります。
10 行目の \code{assert} は、条件が成り立たなければ、その場でプログラムを止める文です。ここで形を確かめます。
正解 \code{y} は、$N$ 個の番号が並んだ形 $(N,)$ のはずです。
\code{f"..."} は、波括弧の中の値を埋め込んだ文字列で、止まったときに表示されます。
11 行目で、先ほどの \code{softmax} を呼び、確率 \code{p} に直します。
要は、12 行目です。\code{torch.arange(N)} は、$(0,1,\dots,N-1)$ という番号のテンソルです。
角括弧に、行の番号のテンソルと、列の番号のテンソルを並べて渡すと、\textbf{同じ順番どうしを組にした位置の要素}が取り出されます（\figref{fig:pt-pick}）。
$N=2$、\code{y} が $(0,2)$ なら、\code{p[0, 0]} と \code{p[1, 2]} の 2 個です。
範囲を二つ並べたスライス（\figref{fig:pt-indexing}の右）が「かたまり」を取り出したのとは、動きが違います。
各行から、正解の列の確率だけを抜き出したことになります。式\eqref{eq:pt-ce}の $p_{n,y_n}$ です。
13 行目で、対数を取り、符号を変え、平均します。式\eqref{eq:pt-ce}の $\sum$ が、\code{for} 文なしで書けました。
\figref{fig:pt-pick}の例を、最後まで手で計算してみます。取り出した確率は 0.7 と 0.6 です。
$-\log 0.7\approx 0.357$、$-\log 0.6\approx 0.511$ なので、損失は $(0.357+0.511)\div 2\approx 0.434$ です。
\figref{fig:pt-ce-flow}に、形の流れをまとめました。コードを読むときは、このように、形がどう変わるかを 1 行ずつ追ってください。

\begin{figure}[tbp]
\centering
\begin{tikzpicture}[x=11mm,y=5.4mm,font=\footnotesize]
  \fill[figorange] (0,1) rectangle (1,2); \fill[figorange] (2,0) rectangle (3,1);
  \draw[black!45,step=1] (0,0) grid (3,2);
  \foreach \v/\x/\y in {0.7/0/1,0.2/1/1,0.1/2/1,0.1/0/0,0.3/1/0,0.6/2/0} \node at (\x+0.5,\y+0.5) {\v};
  \foreach \c in {0,1,2} \node[lbl] at (\c+0.5,2.4) {列 \c};
  \node[lbl,anchor=east] at (-0.1,1.5) {行 0}; \node[lbl,anchor=east] at (-0.1,0.5) {行 1};
  \node[anchor=east] at (-0.9,1) {\texttt{p}\ \ 形 $(2,3)$};
  \node[anchor=west,align=left] at (3.4,1.1) {行の番号\ \ \texttt{torch.arange(2)}\ $=(0,1)$\\列の番号\ \ \texttt{y}\ $=(0,2)$\\[0.5mm]\texttt{p[torch.arange(2), y]}\ $=(0.7,\ 0.6)$};
\end{tikzpicture}
\caption{番号のテンソルを二つ並べたインデックス。1 個目どうし、2 個目どうしを組にした位置（行 0・列 0 と、行 1・列 2）の要素が取り出される。}
\label{fig:pt-pick}
\end{figure}

\begin{figure}[tbp]
\centering
\begin{tikzpicture}[x=1mm,y=1mm,font=\footnotesize,
  s/.style={box,minimum width=16mm,minimum height=9mm,align=center,inner sep=1pt},
  o/.style={font=\scriptsize,align=center,text=bookblue}]
  \node[s] (z) at (0,0) {\texttt{z}\\$(N,V)$};
  \node[s] (p) at (27,0) {\texttt{p}\\$(N,V)$};
  \node[s,fill=figorange,draw=bookorange] (pc) at (60,0) {正解の確率\\$(N,)$};
  \node[s,fill=figorange,draw=bookorange] (l) at (88,0) {$-\log$\\$(N,)$};
  \node[s,fill=figred,draw=bookaccent] (m) at (113,0) {損失\\数 1 個};
  \draw[arr] (z) -- (p); \draw[arr] (p) -- (pc); \draw[arr] (pc) -- (l); \draw[arr] (l) -- (m);
  \node[o] at (13.5,8) {softmax};
  \node[o] at (43.5,9.5) {正解の列を取り出す\\\texttt{p[arange(N), y]}};
  \node[o] at (74,8) {\texttt{-log}};
  \node[o] at (100.5,8) {\texttt{mean}};
  \node[lbl,align=center] at (0,-10.5) {$(2,4096)$\\ロジット};
  \node[lbl,align=center] at (27,-10.5) {$(2,4096)$\\各行の和は 1};
  \node[lbl,align=center] at (60,-10.5) {$(0.0019,\ 0.5590)$\\正解の確率};
  \node[lbl,align=center] at (88,-10.5) {$(6.25,\ 0.58)$\\予測ごとの損失};
  \node[lbl,align=center] at (113,-10.5) {$3.416$\\平均};
\end{tikzpicture}
\caption{交差エントロピーの計算での、形の流れ。下の段は、基盤モデルに「どうも仕方がない。」を与えたときの実際の値である。}
\label{fig:pt-ce-flow}
\end{figure}

PyTorch には、同じ計算をする \code{torch.softmax} と \code{F.cross_entropy} が用意されています（\code{F} は \code{torch.nn.functional} の習慣的な別名です）。
乱数のロジットで比べると、自作の関数の結果は、これらと 6 桁目まで一致しました（\file{ch02/softmax_ce.py}）。
損失は、自作が 4.854022\dots、\code{F.cross_entropy} が 4.854023\dots です。
小数の計算には丸めの誤差があるので、二つのテンソルを比べるときは、\code{==} ではなく \code{torch.allclose(a, b)} を使います。
「ごく小さな差を除いて、全部の要素が等しいか」を、\code{True} か \code{False} で返す関数です。softmax どうしの比較は \code{True} でした。
では、本物の基盤モデルで試します。
「どうも仕方がない。」という文を与え、モデルの予測の損失を、自作の関数で計算します。

\snippet{基盤モデルの損失を、自作の関数で再現する}{code:pt-ce-real}{ch02/softmax_ce.py}{real}

\outputfile{ch02/out/softmax_ce_3.txt}

1〜2 行目で、トークナイザ（文をトークンの番号の列に直す道具。\chapref{ch:tokenizer}で作ります）と、基盤モデルを読み込みます。
3〜4 行目で、文を番号の列に直して表示します。
\code{tok.encode} は文を番号のリストに直し、\code{tok.decode} は番号のリストを文字に戻します。
4 行目の \code{[tok.decode([i]) for i in ids]} は、リスト内包表記という Python の書き方です。\code{ids} の番号を 1 個ずつ文字に戻し、結果をリストに並べます。
出力の 1 行目のとおり、この文は「どうも」「仕方が」「ない。」の 3 トークンです。
5 行目で、入力 \code{x} と正解 \code{y} を作ります。
\code{ids[:-1]} は最後を除いた「どうも」「仕方が」、\code{ids[1:]} は先頭を除いた「仕方が」「ない。」です。
モデルは、入力のどの位置でも、「そこまで読んだ次のトークン」を予測します。
「どうも」の次は「仕方が」、「どうも仕方が」の次は「ない。」という、2 個の予測問題になります。
6〜7 行目で、モデルに予測させます。正解も一緒に渡すと、モデルが自分で計算した損失 \code{loss} も返ってきます。
6 行目の \code{with torch.no_grad():} は、\secref{sec:pt-autograd}の二つ目の決まりです。予測するだけなので、計算の記録を取りません。
モデルは、文の束を形（文の数, トークンの数）で受け取る約束なので、\code{x[None]} で形 $(1,2)$ にして渡します。
\code{x[None]} は、先頭に長さ 1 の軸を足す書き方です（\tabref{tab:pt-shape-ops}）。形 $(2,)$ が $(1,2)$ になり、「2 トークンの文が 1 個」という束になります。
返ってくる \code{logits} の形は $(1,2,4096)$ で、8 行目の \code{logits[0]} で 1 個目の文を取り出すと、形 $(2,4096)$ になります。

ここからが、自作の関数の出番です。
9 行目で softmax を掛け、10 行目で、正解のトークンに割り当てられた確率を取り出します。結果は 0.0019 と 0.5590 でした。
「どうも」のあとには、さまざまな言葉が続きうるので、次が「仕方が」である確率は 0.19\% しかありません。
一方、「どうも仕方が」の次が「ない。」である確率は 55.9\% です。\chapref{ch:intro}と\chapref{ch:math}で見た「約 56\%」です。
11 行目では、2 個目の予測 \code{p[1]} で確率が最大の番号を \code{argmax} で求め、文字に戻しています。確かに「ない。」です。

12 行目が、この章の到達点です。
自作の \code{cross_entropy} の値と、モデルが内部で計算した損失は、どちらも 3.416158\dots で、\textbf{表示された全桁が一致しました}。
260 万個のパラメータを持つモデルの「損失」は、この章で書いた数行の計算と、同じものだったのです。
内訳は、$-\log 0.00193\approx 6.250$ と $-\log 0.5590\approx 0.582$ の平均で、$(6.250+0.582)\div 2=3.416$ です（0.0019 は、丸める前は 0.00193 です）。
この値を、何も知らないモデルと比べます。
何も知らないモデルは、4,096 個の候補のどれにも、同じ確率 $1/4096$ を付けます。
そのときの損失は $-\log(1/4096)=\log 4096\approx 8.318$ でした（\chapref{ch:math}）。
8.318 から 3.416 へ、学習によって損失が下がったことが分かります。
さらに、\code{torch.no_grad()} を外して計算し直し、損失に \code{backward()} を呼ぶと、36 個のテンソルすべてに、同じ形の勾配が付きました（\file{ch02/softmax_ce.py}）。
260 万個の数の一つひとつについて、「どちらに動かせば損失が減るか」が分かったのです。学習の材料は、これでそろいました。

最後に、実際の LLM と規模を比べます（\tabref{tab:pt-scale}）。この章で扱ったテンソルは、大きくても数百万要素でした。

\begin{table}[tbp]
\centering
\caption{テンソルの大きさの比較。「トークンの表」は、各トークンに $d$ 個の数を割り当てる、形 $(V,d)$ の行列である。
GPT-2 と Llama 3.1 の値は、公開されている構成による\cite{radford2019language,dubey2024llama3}。}
\label{tab:pt-scale}
\small
\begin{tabular}{@{}lrrr@{}}
\toprule
 & 本書のミニGPT & GPT-2（小） & Llama 3.1 405B \\
\midrule
トークンの表の形 $(V,d)$ & $(4096,\ 192)$ & $(50257,\ 768)$ & $(128256,\ 16384)$ \\
その要素数 & 約 79 万 & 約 3,860 万 & 約 21 億 \\
パラメータの総数 & 約 258 万 & 約 1.24 億 & 約 4,050 億 \\
\texttt{float32} で持つときのメモリ & 約 10 MB & 約 0.5 GB & 約 1.6 TB \\
\bottomrule
\end{tabular}
\end{table}

Llama 3.1 405B（405B は、パラメータが 4,050 億個という意味です）は、テンソル 1 個だけで約 21 億要素あり、本書のモデル全体の約 800 倍です（約 21 億 $\div$ 約 258 万）。
メモリは、パラメータの総数に 4 バイトを掛けたものです。
4,050 億 $\times$ 4 バイトは 1 兆 6,200 億バイトで、約 1.6 テラバイトになり、PC のメモリには載りません。
そのため、多数の GPU にテンソルを分けて置き、1 要素を 16 ビットで持つ型を使います（\chapref{ch:gpu}）。
しかし、違うのは\textbf{形の中の数字だけ}です。
形を確かめ、\code{@} で行列積を取り、ブロードキャストで形を合わせ、\code{backward()} で勾配を求める。この章の道具は、どの規模でもそのまま使われています。

その道具を、あとから引けるように、\tabref{tab:pt-cheatsheet}にまとめました。
後の章のコードを読んでいて書き方を忘れたら、この表から、該当する節に戻ってください。

\begin{table}[tbp]
\centering
\caption{この章の道具の早見表。左の列から「やりたいこと」を探し、右端の節に戻る。\texttt{S}、\texttt{A}、\texttt{x} などは、それぞれの節の例のテンソルである。}
\label{tab:pt-cheatsheet}
\footnotesize
\begin{tabular}{@{}lll@{}}
\toprule
やりたいこと & 書き方 & 戻る場所 \\
\midrule
形・型・要素数を調べる & \texttt{x.shape}、\texttt{x.dtype}、\texttt{x.numel()} & \secref{sec:pt-tensor} \\
行・列・範囲を取り出す & \texttt{A[1]}、\texttt{A[:, 2]}、\texttt{A[0:2, 1:3]} & \secref{sec:pt-tensor} \\
区切り方を変える & \texttt{x.reshape(3, 4)}（\texttt{-1} は自動） & \secref{sec:pt-reshape} \\
行と列を入れ替える & \texttt{A.T} & \secref{sec:pt-reshape} \\
軸をつぶして集計する & \texttt{S.sum(dim=0)}、\texttt{S.mean(dim=1)} & \secref{sec:pt-broadcast} \\
つぶした軸を長さ 1 で残す & \texttt{S.mean(dim=1, keepdim=True)} & \secref{sec:pt-broadcast} \\
最大値の位置、最大値そのもの & \texttt{S.argmax(dim=1)}、\texttt{S.max(dim=1).values} & \secref{sec:pt-broadcast} \\
長さ 1 の軸を取り除く、足す & \texttt{pred.squeeze(1)}、\texttt{x[None]} & \secref{sec:pt-broadcast}、\tabref{tab:pt-shape-ops} \\
多数のデータに $\mW\vx+\vb$ を適用する & \texttt{Y = X @ W.T + b} & \secref{sec:pt-matmul} \\
勾配を求める & \texttt{requires\_grad=True}、\texttt{f.backward()}、\texttt{x.grad} & \secref{sec:pt-autograd} \\
勾配を 0 に戻す & \texttt{x.grad.zero\_()} & \secref{sec:pt-autograd} \\
記録を取らずに計算する & \texttt{with torch.no\_grad():} & \secref{sec:pt-autograd} \\
形の約束を書く & \texttt{assert y.shape == (N,)} & \secref{sec:pt-final} \\
各行から正解の列を取り出す & \texttt{p[torch.arange(N), y]} & \secref{sec:pt-final} \\
\bottomrule
\end{tabular}
\end{table}

この節では、エラー文の読み方と \code{assert} を学び、自作の softmax と交差エントロピーが、本物のモデルの損失を再現することを確かめました。

\begin{summary}
\item \textbf{配列演算}: 数を 1 個ずつ回る \code{for} 文は、配列演算より 1,000 倍以上遅い。
\item \textbf{形と型}: テンソルは、1 列に並んだ数と、形と、型からなる。モデルは、その集まりである。
\item \textbf{形を変える}: \code{reshape} は区切り方を変え、転置は軸を入れ替える。形が同じでも、中身は違いうる。
\item \textbf{集計とブロードキャスト}: \code{dim} は「つぶす軸」を指定する。形の違う相手とは、右端でそろえ、長さ 1 の軸を引き伸ばして計算する。
\item \textbf{1 データ 1 行}: 数式の $\vy=\mW\vx+\vb$ は、コードでは \code{Y = X @ W.T + b} と書く。
\item \textbf{自動微分}: 印を付けて計算し、\code{backward()} を呼ぶと、勾配が \code{.grad} に入る。
\item \textbf{形を見よ}: 形の取り違えは、エラーにならないことがある。\code{.shape} を確かめ、\code{assert} で形の約束を書く。
\end{summary}

\section*{次の章へ}

道具はそろいました。テンソルで数式を計算でき、損失を測ることができ、その勾配も 1 行で求められます。
しかし、この章では、モデルの数を 1 個も書き換えていません。
基盤モデルの損失 3.416 と、260 万個の数それぞれを「どちらに動かせば損失が減るか」は、手に入りました。
それなら、どれだけ動かせばよいのでしょうか。大きく動かすほど、速く損失が下がるのでしょうか。
損失を下げていく「学習」とは、具体的に何をすることなのでしょうか。
\chapref{ch:ml}では、最も単純なモデルを使って、損失を測り、勾配を求め、数を書き換えるという繰り返しを、自分の手で動かします。
その中で、データ・モデル・損失・最適化という、機械学習の枠組みを組み立てます。

\section*{さらに学ぶために}

\begin{itemize}
\item \cite{paszke2019pytorch}：PyTorch の設計を、開発者たち自身が説明した論文です。
\item \cite{zhang2023d2l}：第 2 章が、本章と同じ範囲を、実行できるノートブックで扱っています。
\item \cite{karpathy2022zerotohero}：動画の講義です。\code{keepdim} の落とし穴を、実演で見られます。
\item \cite{baydin2018automatic}：自動微分の解説論文です。\chapref{ch:nn}のあとに読むと、よく分かります。
\end{itemize}

\begin{exercises}
\exfirst{1〜3}
\item \textbf{形を当てる}。形 $(2,3,4)$ のテンソルに、形 $(4,)$、$(3,1)$、$(2,1,1)$、$(3,)$、$(2,4)$ のテンソルをそれぞれ足します。結果の形（またはエラー）を規則から予想し、実行して確かめてください。
\item \textbf{割合に直す}。回数を数えた 3 行 3 列の表 \code{C} の各行を、その行の合計で割りたいとします。\code{C / C.sum(dim=1)} は、エラーにならずに間違った答えを返します。何が何で割られたかを説明し、直してください（答え合わせは \file{ch02/broadcast_bug.py}）。
\item \textbf{自動微分の検算}。$f(\vx)=\sum_{i}x_i^{3}$ の、$\vx=(1,2,3)$ での勾配を自動微分で求め、手計算と比べてください。\code{x.grad} を 0 に戻さずにもう一度計算すると、どうなるでしょうか。
\item \textbf{あふれない softmax}。ロジット \code{torch.tensor([[1000.0, 999.0, 998.0]])} を\coderef{code:pt-softmax}の \code{softmax} に渡すと、何が返るでしょうか。各行から最大値を引いて計算するように直し、\code{torch.softmax} と比べてください。ヒント: 行ごとの最大値は \code{z.max(dim=-1, keepdim=True).values} です。\code{nan} は「数でない値」です。
\item \textbf{損失の勾配を観察する}。予測 1 個について、\code{cross_entropy} の値をロジットで微分してください。ロジットは形 $(1,3)$ の \code{[[2.0, 1.0, 0.0]]}、正解は \code{torch.tensor([0])} とします。勾配と softmax の出力 $\vp$ には、どんな関係があるでしょうか（\chapref{ch:ml}で導きます）。
\end{exercises}
